\chapter{衍射积分、孔径与光栅分辨率}
\section{电磁波、复数平均与光强}\label{th:o3wave}
在均匀、线性、各向同性、无自由电荷电流的介质中，$\vct D=\varepsilon\vct E$、$\vct B=\mu\vct H$，麦克斯韦方程给$\nabla\cdot\vct E=0$。对法拉第定律取旋度，利用$\nabla\times\nabla\times\vct E=\nabla(\nabla\cdot\vct E)-\nabla^2\vct E$，得到
\[
 \nabla^2\vct E-\varepsilon\mu\partial_t^2\vct E=0.
\]
代入$\vct E=\vct E_0\ee^{i(\vct k\cdot\vct r-\omega t)}$，得到$\vct k\cdot\vct E_0=0$、$k^2=\varepsilon\mu\omega^2$。相速度$v=1/\sqrt{\varepsilon\mu}=c/n$，磁场$\vct B_0=\vct k\times\vct E_0/\omega$，所以$\vct E,\vct H,\vct k$构成定向正交关系。若介质吸收或各向异性，这些实数和共线关系需重新讨论。

\subsection{时间平均中的共轭}\label{th:o3intensity}
同频复振幅$A,B$的实部乘积平均为$\langle\operatorname{Re}(A\ee^{-i\omega t})\operatorname{Re}(B\ee^{-i\omega t})\rangle=\tfrac12\operatorname{Re}(AB^*)$。可把每个实部写成两个指数之和，乘积中的$\ee^{\pm2i\omega t}$在整周期平均中消失，只留下互为共轭的两项。

因而平均坡印廷矢量$\langle\vct S\rangle=\tfrac12\operatorname{Re}(\vct E_0\times\vct H_0^*)$。无损平面波的光强为$I=E_0^2/(2\mu v)$。非磁性介质$\mu=\mu_0$时，$I=n\varepsilon_0cE_0^2/2$。这个强度沿传播方向定义；穿过倾斜表面的法向能流还要乘余弦。

复振幅叠加先做矢量求和，再作时间平均。若两同偏振场为$A_1\ee^{i\phi_1}$和$A_2\ee^{i\phi_2}$，合相位应由实虚部使用$\operatorname{atan2}$确定，不能把各自相位直接平均。只有振幅相同、相位差在适当支路内，角平分线解释才可直接使用。

\section{从惠更斯积分到远场傅里叶积分}\label{th:o3integral}
设孔径面为$z=0$，孔内复振幅为$U(Q)$。标量近似下，观察点$P$的场可写为
\[
 U(P)\simeq\frac1{i\lambda}\iint_\Sigma
 U(Q)\frac{\ee^{ikr}}r\frac{\cos\theta+\cos\theta_0}2\dd S,
\]
其中$r=|P-Q|$，$\theta,\theta_0$分别描述出射和入射方向与孔面法向的夹角。该式忽略孔边附近的严格矢量边界场，适于孔径远大于波长、观察方向非极端掠射的情形。

观察屏在$z=L$，孔径坐标$(\xi,\eta)$、观察坐标$(x,y)$。近轴距离展开
\[
 r\simeq L+\frac{(x-\xi)^2+(y-\eta)^2}{2L}.
\]
幅度因子$1/r$可近似$1/L$，但指数中的小距离会乘大波数，不能同时全部省略。于是
\begin{align}\label{eq:o3fresnel}
 U(x,y)&\simeq\frac{\ee^{ikL}}{i\lambda L}
 \ee^{ik(x^2+y^2)/(2L)}\\
 &\quad\times\iint U(\xi,\eta)
 \ee^{ik(\xi^2+\eta^2)/(2L)}
 \ee^{-ik(x\xi+y\eta)/L}\dd\xi\dd\eta.\notag
\end{align}
保留孔径二次相位是菲涅耳衍射。若孔径特征半宽$a$满足$a^2/(\lambda L)\ll1$，二次相位变化小，可略去，得到夫琅禾费积分。薄透镜后焦面在近轴模型中也可给出孔径的傅里叶分布，无须真的把屏移到无限远；其相位补偿已纳入透镜传播。

\Diagram{fig:o3integral}{孔径点$Q$到屏上点$P$的路径及坐标}{\draw[curve](0,-1.6)--(0,1.7);\draw[curve](5,-1.6)--(5,2.2);\draw[guide](0,0)--(5,0);\fill(0,.6)circle(1.4pt)node[left]{$Q(\xi,\eta)$};\fill(5,1.4)circle(1.4pt)node[right]{$P(x,y)$};\draw[curve](0,.6)--(5,1.4)node[midway,above]{$r$};\node[below]at(2.5,0){$L$};}

\section{单缝的积分与相量圆弧}\label{th:o3slit}
均匀照亮宽度$a$的长直缝，先只处理宽度方向。沿$x$方向观察角为$\theta$，孔内坐标$\xi\in[-a/2,a/2]$，忽略共同相位因子，
\[
 U(\theta)=C\int_{-a/2}^{a/2}\ee^{-ik\xi\sin\theta}\dd\xi
 =C\left[\frac{\ee^{-ik\xi\sin\theta}}{-ik\sin\theta}\right]_{-a/2}^{a/2}
 =Ca\frac{\sin\beta}{\beta},\quad \beta=\frac{\pi a\sin\theta}\lambda.
\]
因此$I(\theta)=I(0)(\sin\beta/\beta)^2$。$\beta=0$取极限1，不能因分子为零判作暗纹。暗纹$a\sin\theta_m=m\lambda$，$m=\pm1,\pm2,\ldots$。中央主瓣首零全角宽近轴为$2\lambda/a$。

若把缝坐标改为$[0,a]$，合场多一个$\ee^{-i\beta}$因子，但强度不变；对后续与其他孔径的相干叠加，这个位置相位则不能删掉。场振幅的实值sinc随旁瓣交替变号，强度始终非负。画图时需说明纵轴是场、场模还是强度。

\begin{example}\label{ex:o3phasor}
把单缝分为无穷多个等宽小元，以相量首尾相接求合场，解释第一个暗纹和第一旁瓣场的反号。
\Diagram{fig:o3phasorq}{缝上各小元到同一远场方向的相位线性变化}{\draw[curve](0,-1.6)--(0,1.6);\foreach\y in{-1.4,-.7,0,.7,1.4}{\draw[->](0,\y)--(3,{\y+.8});}\node[left]at(0,0){$a$};\node[right]at(3,.8){$\theta$};}
\end{example}
\begin{solution}
图\ref{fig:o3phasorq}中相邻小元振幅相等、相位递增相同，首尾相接形成圆弧。整缝相位变化为$ka\sin\theta=2\beta$，弧长与轴上场$Ca$成正比。若相量圆半径为$R_p$，弧长$2R_p\beta=Ca$，合场弦长为$2R_p\sin\beta=Ca\sin\beta/\beta$。
\Diagram{fig:o3phasora}{相量圆弧的端点连线是合场；一整圆对应首零}{\draw[curve](0,0)circle(1.4);\draw[->,thick](1.4,0)arc(0:150:1.4);\draw[->,very thick](1.4,0)--(-1.212,.7);\draw[guide](0,0)--(1.4,0);\draw[guide](0,0)--(-1.212,.7);\node[above]at(0,1.7){圆心角$2\beta$};\node[below]at(0,-1.7){整圆时端点重合};}
图\ref{fig:o3phasora}的一整圆对应$2\beta=2\pi$，相量闭合为零。超过整圆后，合弦相对原公共相位反号，故第一旁瓣的实场sinc为负；模平方仍为正。上下半缝的面元配对可解释首零，圆弧构造还能给出任意角度的振幅与相位。
\end{solution}

\begin{example}\label{ex:o3rectangle}
矩形孔径宽$a$、高$H$，均匀平面波正入射。求远场二维强度，并讨论长缝近似的适用方向。
\Diagram{fig:o3rectangleq}{以孔中心为坐标原点的矩形孔径}{
\draw[axis](-1.8,0)--(2,0)node[right]{$\xi$};
\draw[axis](0,-2.1)--(0,2.3)node[above]{$\eta$};
\draw[curve](-1.1,-1.6)rectangle(1.1,1.6);
\draw[<->](-1.1,-1.9)--(1.1,-1.9)node[midway,below]{$a$};
\draw[<->](1.5,-1.6)--(1.5,1.6);\node[right]at(1.5,.7){$H$};
\node[below left]at(0,0){$O$};
}
\end{example}
\begin{solution}
图\ref{fig:o3rectangleq}中取孔中心为原点，$-a/2<\xi<a/2$、$-H/2<\eta<H/2$。观察方向余弦为$l_x,l_y$，令$q_x=kl_x$、$q_y=kl_y$，场积分为
\[
 U=C\int_{-a/2}^{a/2}\ee^{-iq_x\xi}\dd\xi
 \int_{-H/2}^{H/2}\ee^{-iq_y\eta}\dd\eta.
\]
两个积分分别取原函数$\ee^{-iq_x\xi}/(-iq_x)$、$\ee^{-iq_y\eta}/(-iq_y)$，在$\pm a/2$和$\pm H/2$代入，得到$2\sin(q_xa/2)/q_x$和$2\sin(q_yH/2)/q_y$；零横向波矢处直接积分为$a,H$。于是
\[
 U=CaH\operatorname{sinc}(\pi a l_x/\lambda)
 \operatorname{sinc}(\pi H l_y/\lambda),
\]
其中$\operatorname{sinc}q=\sin q/q$，零点取连续极限1。
\Diagram{fig:o3rectanglea}{孔径越长的方向，首零边界越靠近中心}{\draw[axis](-2.4,0)--(2.6,0)node[right]{$l_x$};\draw[axis](0,-1.5)--(0,1.7)node[above]{$l_y$};\draw[curve](-1.8,-.45)rectangle(1.8,.45);\node[above]at(0,1){首零边界};\node[below]at(1.8,-.5){$\lambda/a$};\node[left]at(0,.6){$\lambda/H$};}
图\ref{fig:o3rectanglea}中$H>a$，所以沿$y$的角分布更窄：首零方向余弦分别为$|l_x|=\lambda/a$和$|l_y|=\lambda/H$。强度为两个sinc平方之积，中央矩形亮斑在宽孔对应方向上较窄。取$l_y=0$切面，第二因子等于1，才恢复一维长缝结果。

轴上场正比孔面积$aH$，轴上强度正比$(aH)^2$；但两方向角尺度分别含$1/a$和$1/H$，对整个远场积分后总功率仍只正比$aH$。若仅增大$H$，$x$切面的归一化形状不变，绝对峰强和$y$方向宽度却都改变。孔高不能在二维问题中被无条件省去。
\end{solution}

\section{双缝包络与有限宽光栅}\label{th:o3grating}
两条宽$a$的缝，其中心相隔$b$。利用平移对应相位因子，合场为$U=U_a[1+\ee^{-ikb\sin\theta}]$，所以
\[
 I=4I_a(0)\operatorname{sinc}^2\beta\cos^2\alpha,
 \quad\beta=\pi a\sin\theta/\lambda,\quad
 \alpha=\pi b\sin\theta/\lambda.
\]
对$N$条等宽等间距缝，同理
\begin{equation}\label{eq:o3grating}
 I=I_a(0)\operatorname{sinc}^2\beta
 \left[\frac{\sin N\alpha}{\sin\alpha}\right]^2.
\end{equation}
单缝包络决定每一级主峰获得多少光能，多缝因子决定主峰宽度和细结构。改变$a/b$会改变缺级和效率，改变被照亮缝数$N$才直接改变同级主峰宽度。这里所有缝照明振幅相同、相位线性递增；照明不均匀时须改成带权和。

\begin{example}\label{ex:o3missing}
双缝或光栅的周期$b=3a$。按干涉因子的主极大条件，求单缝中央包络内的级次，指出缺级，并画出主极大与包络的对应位置。另讨论传播条件$|\sin\theta|\le1$对这些级次的限制。
\Diagram{fig:o3missingq}{有限缝宽$a$和中心周期$b=3a$}{\draw[curve](0,-.5)rectangle(.5,1);\draw[curve](1.5,-.5)rectangle(2,1);\draw[<->](0,-.9)--(.5,-.9)node[midway,below]{$a$};\draw[<->](.25,1.4)--(1.75,1.4)node[midway,above]{$b$};}
\end{example}
\begin{solution}
干涉因子的主极大满足$b\sin\theta=m\lambda$，单缝零点满足$a\sin\theta=p\lambda$，$p\ne0$。图\ref{fig:o3missingq}取$b=3a$，两条件重合时$m=3p$，故在能够传播的级次中，$\pm3,\pm6,\ldots$缺级。
\Diagram{fig:o3missinga}{以$b\sin\theta/\lambda$为横轴，包络首零在$\pm3$}{\draw[axis](-3.7,0)--(3.9,0)node[right]{$m$};\draw[axis](0,0)--(0,2.7)node[above]{强度};\draw[curve,domain=.01:3.5,samples=70]plot(\x,{2.2*(sin(60*\x)/(pi*\x/3))^2});\draw[curve,domain=-3.5:-.01,samples=70]plot(\x,{2.2*(sin(60*\x)/(pi*\x/3))^2});\foreach\x in{-2,-1,1,2}{\draw[very thick](\x,0)--(\x,{2.2*(sin(60*\x)/(pi*\x/3))^2});}\draw[very thick](0,0)--(0,2.2);\node[below]at(3,0){$3$};\node[below]at(-3,0){$-3$};}
单缝中央包络严格内部为$|a\sin\theta/\lambda|<1$，代入主极大条件得到$|m|<3$，候选级次为$m=0,\pm1,\pm2$。还须逐个检验传播条件
\[
 |\sin\theta_m|=\frac{|m|\lambda}{b}\le1.
\]
当$b>2\lambda$时，五个候选级次都在非掠射方向传播；当$\lambda<b<2\lambda$时只剩$0,\pm1$；当$0<b<\lambda$时只剩零级。$b=\lambda$或$2\lambda$时，相应一级或二级位于掠射方向，是否被接收还取决于观察几何。一般应取整数集合$\{m:|m|<3,\ |m|\lambda/b\le1\}$。

图\ref{fig:o3missinga}的横轴是无量纲方向变量$b\sin\theta/\lambda$；图示范围需取$b>3\lambda$，才能同时画出两侧首零。竖线标干涉因子的主极大位置，平滑线标单缝包络；实际有限$N$的峰有有限宽。总强度是两因子的乘积，除对称中心等特殊位置外，其局部峰还可能因包络斜率而略有偏移。
\end{solution}

\section{分辨判据与孔径宽度}
\begin{example}\label{ex:o3sodium}
要在第$m$级分辨波长589.0 nm和589.6 nm两线，按瑞利判据至少需照亮多少条缝？若只照亮200条缝，最低可用级次为何？
\Diagram{fig:o3sodiumq}{两条谱线在同一级产生相邻主峰}{\draw[axis](0,0)--(5.5,0)node[right]{$\theta$};\draw[curve](1,0)..controls(1.5,2.5)and(1.8,2.5)..(2.4,0);\draw[curve](2,0)..controls(2.5,2.5)and(2.8,2.5)..(3.4,0);\node[above]at(1.7,2){589.0};\node[above]at(2.7,2.7){589.6};}
\end{example}
\begin{solution}
取平均波长$\bar\lambda=589.3$ nm、差$\Delta\lambda=.6$ nm，所需分辨本领约982.17。图\ref{fig:o3sodiumq}按瑞利比较要求$mN\geq982.17$，因此$N_{\min}=\lceil982.17/m\rceil$。
\Diagram{fig:o3sodiuma}{提高级次缩小所需缝数，但必须有该传播级次}{\draw[axis](0,0)--(4.8,0)node[right]{$m$};\draw[axis](0,0)--(0,3.2)node[above]{$N_{\min}$};\foreach\x/\y/\t in{1/2.95/983,2/1.476/492,3/.984/328,4/.738/246}{\fill(\x,\y)circle(1.5pt);\node[right]at(\x,\y){$\t$};}}
一级至少983条，二级492条。若$N=200$，则$m\geq\lceil982.17/200\rceil=5$。图\ref{fig:o3sodiuma}的提高级次并非无条件可行，还须$|m\lambda/b|\leq1$、该级不缺级且效率足够。只满足$mN$而传播角不存在时，不能实现分辨。
\end{solution}

\begin{example}\label{ex:o3sparrow}
两条等强、互不相干谱线产生相同的窄主峰，单线近似$F(x)=(\sin x/x)^2$。两峰中心为$\pm a$。以总曲线中点由极大转为极小作为双峰边界，推导所需间距，并与瑞利间距比较。
\Diagram{fig:o3sparrowq}{两条单线的强度相加，峰中心相距$2a$}{\draw[axis](-3,0)--(3.4,0)node[right]{$x$};\draw[curve,domain=-2.7:2.7,samples=90]plot(\x,{1.5*exp(-2*(\x-1)^2)});\draw[guide,domain=-2.7:2.7,samples=90]plot(\x,{1.5*exp(-2*(\x+1)^2)});\node[below]at(-1,0){$-a$};\node[below]at(1,0){$a$};}
\end{example}
\begin{solution}
图\ref{fig:o3sparrowq}只示意两中心位置，定量使用sinc平方。总强度$G(x)=F(x-a)+F(x+a)$为偶函数，$G'(0)=0$、$G''(0)=2F''(a)$。曲率边界为$F''(a)=0$。完整求导给
\[
 F'(x)=\frac{\sin2x}{x^2}-\frac{2\sin^2x}{x^3},
\]
\[
 F''(x)=\frac{2\cos2x}{x^2}-\frac{4\sin2x}{x^3}
 +\frac{6\sin^2x}{x^4}.
\]
在$x=a$处通分并用$2\sin^2a=1-\cos2a$，得
\[
 F''(a)=\frac{(2a^2-3)\cos2a-4a\sin2a+3}{a^4}.
\]
取第一非零正根，数值$a\simeq1.30308171$，所以两中心间距$2a\simeq2.60616$。
\Diagram{fig:o3sparrowa}{曲率判据的临界中点变平，超过临界后出现中央凹陷}{\draw[axis](-3,0)--(3.3,0)node[right]{$x$};\draw[axis](0,0)--(0,2.6)node[above]{$G$};\draw[curve,domain=-2.8:2.8,samples=120]plot(\x,{(sin((\x-1.30308) r)/(\x-1.30308))^2+(sin((\x+1.30308) r)/(\x+1.30308))^2});\draw[guide](-.6,1.096)--(.6,1.096);}
瑞利判据的中心间距是单峰中心到首零的$\pi$，故图\ref{fig:o3sparrowa}对应的曲率间距约为瑞利间距的$2.60616/\pi\simeq0.82957$。求导时$x^{-2}$的变化与$\sin^2x$的变化同阶，不能把分母固定后仅对分子求导。此判据还依赖两线等强、无噪声和相同形状；实际能否可靠辨认需结合测量条件。
\end{solution}

\begin{example}\label{ex:o3width}
光栅照明总宽度$W$，周期$b$，正常入射。证明任意可传播级次的瑞利分辨本领受$W/\lambda$限制，并解释为何无限缩小周期不能无限提高分辨率。
\Diagram{fig:o3widthq}{照明宽度决定实际参与的缝数}{\draw[curve](0,-.5)rectangle(5,.7);\foreach\x in{.2,.6,1,1.4,1.8,2.2,2.6,3,3.4,3.8,4.2,4.6}{\draw[curve](\x,-.5)--(\x,.7);}\draw[<->](0,-1)--(5,-1)node[midway,below]{$W\simeq Nb$};}
\end{example}
\begin{solution}
图\ref{fig:o3widthq}中$N\simeq W/b$。主极大$b\sin\theta=m\lambda$要求$|m|\leq b/\lambda$。相乘得到$|m|N\lesssim W/\lambda$。
\Diagram{fig:o3widtha}{周期越小虽增加缝数，却降低最高可用级次}{\draw[axis](0,0)--(4.6,0)node[right]{$b/\lambda$};\draw[axis](0,0)--(0,3)node[above]{最高$|m|$};\draw[very thick](0,0)--(1,0)--(1,.7)--(2,.7)--(2,1.4)--(3,1.4)--(3,2.1)--(4,2.1);\draw[guide](0,0)--(4,2.8);}
图\ref{fig:o3widtha}中的级次是整数，实际最大值为$\lfloor b/\lambda\rfloor$。当$b<\lambda$只剩零级，不能靠通常角色散分辨不同波长。上限在掠射方向逼近，近轴角宽公式及实际效率此时需要更细处理；这一上限由实际参与相干叠加的孔径宽度确定。
\end{solution}

\section{闪耀反射光栅}\label{th:o3blaze}
反射光栅每周期的小面可视为有限平面镜。设宏观法线沿$+z$，入射光沿$-z$，小面法线向$+x$倾斜$\alpha$，出射方向相对$+z$成$\theta$。小面切向单位矢量为$(\cos\alpha,-\sin\alpha)$，其上坐标$s$的相位梯度为
\[
 (\vct k_i-\vct k_o)\cdot\vct t
 =k[\sin\alpha-\sin(\theta-\alpha)].
\]
单小面包络在梯度为零时最强，即$\theta=2\alpha$，正是镜面反射方向。周期阵列的主峰仍由$b\sin\theta_m=m\lambda$确定；使某级落在包络峰上，闪耀波长为$\lambda_B=b\sin2\alpha/m$。阴影遮挡、偏振及真实沟槽形状会影响效率，简单标量面元积分只描述主要相位匹配。

\begin{example}\label{ex:o3blaze}
正常入射的闪耀光栅周期$b=2\,\mu\mathrm m$，小面倾角$\alpha=15^\circ$。求二级闪耀波长，并比较此处相邻周期相位和单小面内相位各自的作用。
\Diagram{fig:o3blazeq}{倾斜小面使镜面包络偏向某一衍射级}{\foreach\x in{0,1.5,3,4.5}{\draw[curve](\x,0)--({\x+1.1},-.3)--({\x+1.5},0);}\draw[->](2,2)--(2,-.14);\draw[->](2,-.14)--(3.3,2.11);\node[above]at(2,2){入射};\node[right]at(3.3,2.11){$2\alpha$};}
\end{example}
\begin{solution}
图\ref{fig:o3blazeq}的包络中心$\theta=30^\circ$，二级光栅条件$2\lambda_B=b\sin30^\circ$，所以$\lambda_B=.5\,\mu\mathrm m=500$ nm。
\Diagram{fig:o3blazea}{相同多缝窄峰乘以偏移包络，强度集中在选定级次}{\draw[axis](0,0)--(5.8,0)node[right]{衍射级次};\draw[curve,domain=0:5,samples=80]plot(\x,{2.5*exp(-1.1*(\x-2)^2)});\foreach\x in{0,1,2,3,4,5}{\draw[very thick](\x,0)--(\x,{2.5*exp(-1.1*(\x-2)^2)});}\node[below]at(2,0){$m=2$};}
小面内同相叠加使该角度单缝包络大，周期之间相差$4\pi$使二级主峰相长，图\ref{fig:o3blazea}表示两条件同时满足。闪耀主要提高选定级的能量效率；固定$N,m$下，瑞利分辨本领仍是$mN$，不会因包络更高而自动变成更窄的几何峰。
\end{solution}
\section{衍射积分中的尺度与归一化}
\subsection{几何小量不等于相位小量}
传播距离可展开为$\sqrt{L^2+\rho^2}=L+\rho^2/(2L)-\rho^4/(8L^3)+\cdots$。在幅度因子$1/r$中，只要$\rho/L\ll1$就可近似$r=L$；在相位$kr$中，误差还乘$k=2\pi/\lambda$，要求$k\rho^4/(8L^3)$足够小才能只保留二次项。把三角近轴条件与相位截断条件分开，才能判断菲涅耳近似在特定尺寸下的精度。

夫琅禾费近似还要求孔面二次相位$k(\xi^2+\eta^2)/(2L)$的变化足够小，或已由透镜补偿。增大屏距、缩小孔径都可减小该项；仅让观察角很小还不够。高斯束、直边及大孔径近场计算通常仍须保留二次相位。

在孔径前还有有限距离点光源时，入射到孔面已带一个二次球面相位。传播到屏的二次相位与它相加，等效距离由两段波前曲率共同决定。不能只看孔到屏的$L$来断定远场；若点光源很近，入射曲率项可能仍很大。理想透镜可用反号二次相位补偿相应曲率，形成清楚的傅里叶变换面。

\subsection{远场坐标的量纲}
一维远场场可写$U(x)=C F(q)$，其中$F(q)=\int U_0(\xi)\ee^{-iq\xi}\dd\xi$，$q=kx/L$。$q$量纲为长度倒数，$x$是屏幕长度。若使用空间频率$f_x=x/(\lambda L)$，指数应写$\ee^{-i2\pi f_x\xi}$。这两套约定只差$2\pi$，但单缝零点和高斯宽度的常数必须随之同步。

\Diagram{fig:o3scales}{孔面长度尺度与远场空间频率尺度互为倒数}{\draw[axis](0,1.4)--(5,1.4)node[right]{$\xi$};\draw[very thick](1.5,1.4)--(3.5,1.4);\node[above]at(2.5,1.4){宽$a$};\draw[axis](0,0)--(5,0)node[right]{$q$};\draw[very thick](2,0)--(3,0);\node[below]at(2.5,0){宽度尺度$1/a$};}
图\ref{fig:o3scales}说明孔越宽、角谱越窄，但比较绝对亮度还要保留孔面照明归一化。均匀单缝振幅正比$a$，所以中央强度随$a^2$增长。把每条曲线都除以各自峰值后只剩形状，不能再从归一化图上读取绝对功率变化。

\section{傅里叶积分中的能量守恒}
对平方可积孔场，采用上述无$2\pi$的傅里叶约定，帕塞瓦尔关系为
\[
 \int_{-\infty}^{\infty}|F(q)|^2\dd q
 =2\pi\int_{-\infty}^{\infty}|U_0(\xi)|^2\dd\xi.
\]
它可从双重积分展开$|F|^2$后，使用$\int\ee^{-iq(\xi-\eta)}\dd q=2\pi\delta(\xi-\eta)$得到。自由传播核的系数与屏面变量$\dd x=L\dd q/k$相配合，使积分能流等于孔面透过能流。二维情况沿两个方向各用一次同样关系。

单缝可直接核对功率归一化：孔面振幅恒定时，$F(q)=aU_0\operatorname{sinc}(qa/2)$。换元$u=qa/2$，得$\int|F|^2\dd q=2a|U_0|^2\int\operatorname{sinc}^2u\dd u$，积分正比于$a$。孔宽增大使中央强度升高、角分布变窄，二者共同满足总功率守恒。

\subsection{相位孔与振幅孔}
若孔面仅乘单位模相位$\ee^{i\phi(\xi)}$，总透过功率不变，但远场分布可明显改变。线性相位整体转向，二次相位改变聚焦或观察面，周期相位形成多个衍射级。若孔面乘振幅透过率$|t(\xi)|<1$，则总功率减少；它仍可通过压低边缘权重改善旁瓣，代价是峰宽或通光量。

多束加权的二项式系数和连续高斯照明都是这种思想的离散、连续例子。比较不同孔径时应固定哪一个条件：总输入功率、最大照明强度、有效孔径宽度，还是透过总能量。若约束未说明，峰高、旁瓣与“效率”的比较往往没有唯一答案。

\section{点像、物像与分辨判据}
\subsection{单点响应怎样组合成像}
理想无像差透镜的几何像点在波动理论中仍有有限衍射斑。孔函数的傅里叶变换给出点源的复振幅响应$h(x)$，其强度响应为$|h(x)|^2$。多个彼此不相干物点成像时，各点像强度相加；若物点相干，则应先把复振幅相加。两者即使几何点位完全相同，图样也可能因相对相位不同而改变。

\Diagram{fig:o3pointimages}{不相干物点的点像相加强度，相干物点则先加场}{\fill(0,.7)circle(1.5pt)node[left]{$O_1$};\fill(0,-.7)circle(1.5pt)node[left]{$O_2$};\draw[<->](2,-1.5)--(2,1.5);\draw[curve](4.5,.7)circle(.45);\draw[curve](4.5,-.7)circle(.45);\draw[guide](0,.7)--(4.5,-.7);\draw[guide](0,-.7)--(4.5,.7);}
图\ref{fig:o3pointimages}只标出点像中心与宽度。对两个等强不相干谱线或星点，瑞利判据使用主峰与首暗环重合；对两相干点，交叉项会使中间强度取决于相位，不能沿用同一张强度相加曲线。

\subsection{瑞利、曲率和测量误差}
瑞利判据以主峰与首零点重合作为分离尺度。曲率判据则考察中央由峰变谷的边界，在对称等强模型中令二阶导数为零。对sinc平方线形，后者间距约为瑞利尺度的.8296倍；高斯、圆孔等线形须另求这个系数。

实际谱线很弱时，另一个强峰的翼部与背景噪声可能比中心曲率更重要。高信噪比、已知线形的拟合有时能估计小于传统瑞利间隔的峰位，但这不等于光学点像变窄；估计不确定度取决于光子数、噪声、模型先验和参数相关。竞赛题若指定某种判据，应按该定义计算，不把不同定义混作唯一“最小可分辨距离”。

\subsection{光栅参数中的三个长度}
缝宽$a$、周期$b$、照明总宽$W\simeq Nb$决定不同结构。$a$给单缝包络和缺级；$b$给各级中心角和最高可传播级次；$W$给同级窄峰的宽度。若固定$W$而缩小$b$，虽然$N$增大，可用级次却下降；若固定$b$而增加$N$，主峰中心不动，宽度变窄。

闪耀小面的倾角把单元包络峰移向选定级次，是效率设计；增加总照明宽度才真正增加可参与的相干横向范围。真实光栅还受表面误差、局部周期误差和照明相位不均影响，实际窄峰可能比理想$mN$更宽。这些误差应按具体分布进入孔面场，再求实际线形。

\subsection{孔径、相干照明和点像强度}
夫琅禾费场是孔面复场的傅里叶变换，而非孔面强度的傅里叶变换。若孔面透过率为$t(\xi)$、照明复场为$A(\xi)$，实际参与积分的是$A(\xi)t(\xi)$。均匀照明、振幅加权和附加相位因而改变不同对象：加权改变面元振幅，相位片改变面元方向，遮挡则同时改变通光范围。

对若干点源成像，还要在像面区分相干与不相干叠加。若两个点源经同一系统形成复振幅点响应$h(x-x_1)$、$h(x-x_2)$，完全相干时强度为
\[
 |a_1h(x-x_1)+a_2\ee^{i\phi}h(x-x_2)|^2.
\]
互不相干时交叉项长期平均为零，才得到$|a_1|^2|h(x-x_1)|^2+|a_2|^2|h(x-x_2)|^2$。瑞利判据和本讲等强双峰曲率判据采用后一种模型；在孔径和点源间距固定时，相干性仍会改变中央峰谷。
\Diagram{fig:o3coherentpairs}{实线为同相合场，曲线虚线为反相合场；竖直虚线标出点像中心}{\draw[axis](-.3,0)--(5.7,0)node[right]{$x$};\draw[axis](0,-1)--(0,2.7)node[above]{场};\draw[guide](2,-.4)--(2,2.3);\draw[guide](3.2,-.4)--(3.2,2.3);\draw[curve,domain=.2:5.2,samples=100]plot(\x,{exp(-2*(\x-2)^2)+exp(-2*(\x-3.2)^2)});\draw[dashed,domain=.2:5.2,samples=100]plot(\x,{exp(-2*(\x-2)^2)-exp(-2*(\x-3.2)^2)});\node at(4.7,2){同相};\node at(4.7,-.7){反相};}
图\ref{fig:o3coherentpairs}为平滑实场的示意，实线同相相加，虚线反相相加；它不代替具体圆孔或单缝积分。使用测得的强度曲线讨论分辨时，应先明确各点响应、光源统计和探测平均方式，再选择判据。

远场强度还能检验算式归一化：对完整角域或相应近轴傅里叶平面积分，应与孔面透过功率相符。若只截取主瓣，则所得只是该区域功率；忽略旁瓣后不能继续把主瓣积分称为总功率。若孔径尺寸接近波长、观察角接近掠射或介质边界引入明显矢量效应，标量近轴模型的适用性也须重新判断。

\section{近场边缘与孔径的空间频率}
夫琅禾费近似把孔面场变为空间频率的傅里叶变换；菲涅耳近似则保留每个孔元的二次相位。二者使用同一传播积分，区别在于可省略哪些相位。描述近场时，不能把每条射线的几何阴影边界当成强度突然为零的位置。

直边衍射需要的振荡积分也可由第一讲的复高斯积分确定。先乘收敛因子$\ee^{-\epsilon t^2}$，$\epsilon>0$，在式\eqref{eq:o1gaussintegral}中取$a=\epsilon-i\pi/2$、$b=0$，
\[
 \int_{-\infty}^{\infty}\ee^{-\epsilon t^2+i\pi t^2/2}\dd t
 =\sqrt{\frac{\pi}{\epsilon-i\pi/2}}
 \xrightarrow{\epsilon\to0^+}\sqrt2\ee^{i\pi/4}=1+i.
\]
平方根沿$\operatorname{Re}a>0$连续取支。该收敛极限给出菲涅耳振荡积分；被积函数为偶函数，所以$\int_0^\infty\ee^{i\pi t^2/2}\dd t=(1+i)/2$。实、虚部分别给$C(\infty)=S(\infty)=1/2$，用于下面的无遮挡场归一化。

\begin{example}\label{ex:o3edge}
平面波照射遮住$x<0$的直边，光沿$+z$传播。距边$L$的屏上坐标为$x$，定义$v=x\sqrt{2/(\lambda L)}$。用菲涅耳积分表示归一化强度，求几何边界$x=0$的强度。
\Diagram{fig:o3edgeq}{半平面边缘及几何阴影线}{\draw[curve](0,-1.5)--(0,0);\foreach\y in{-1.4,-1.1,-.8,-.5,-.2}{\draw(0,\y)--(-.3,{\y-.2});}\draw[curve](5,-1.8)--(5,2);\draw[guide](0,0)--(5,0);\draw[->](-1,1)--(1,1);\node[right]at(5,-1){几何阴影};\node[right]at(5,1){照明区};}
\end{example}
\begin{solution}
图\ref{fig:o3edgeq}的透光半平面为$\xi\geq0$。对沿边方向的积分提出共同因子后，垂直边方向场正比$\int_0^\infty\exp[i\pi(x-\xi)^2/(\lambda L)]\dd\xi$。令$t=(x-\xi)\sqrt{2/(\lambda L)}$，$\xi=0$对应$t=v$，$\xi\to\infty$对应$t\to-\infty$，负的微分将上下限反转，故积分正比$\int_{-\infty}^v\ee^{i\pi t^2/2}\dd t$。
\[
 C(v)=\int_0^v\cos(\pi t^2/2)\dd t,\quad
 S(v)=\int_0^v\sin(\pi t^2/2)\dd t.
\]
用$C(\infty)=S(\infty)=1/2$和奇函数性质，分子为$(1/2+C)+i(1/2+S)$。
\Diagram{fig:o3edgea}{按菲涅耳积分计算的直边衍射强度，几何边界处为四分之一}{\draw[axis](-3.2,0)--(3.7,0)node[right]{$v$};\draw[axis](0,0)--(0,3.2)node[above]{$I/I_0$};\draw[curve]plot coordinates {(-3.00000,0.012310) (-2.96425,0.012605) (-2.92849,0.012910) (-2.89274,0.013227) (-2.85698,0.013556) (-2.82123,0.013897) (-2.78547,0.014250) (-2.74972,0.014618) (-2.71397,0.014999) (-2.67821,0.015395) (-2.64246,0.015807) (-2.60670,0.016235) (-2.57095,0.016681) (-2.53520,0.017144) (-2.49944,0.017627) (-2.46369,0.018131) (-2.42793,0.018655) (-2.39218,0.019203) (-2.35642,0.019774) (-2.32067,0.020370) (-2.28492,0.020993) (-2.24916,0.021645) (-2.21341,0.022326) (-2.17765,0.023039) (-2.14190,0.023785) (-2.10615,0.024568) (-2.07039,0.025388) (-2.03464,0.026248) (-1.99888,0.027151) (-1.96313,0.028100) (-1.92737,0.029097) (-1.89162,0.030146) (-1.85587,0.031250) (-1.82011,0.032412) (-1.78436,0.033638) (-1.74860,0.034930) (-1.71285,0.036294) (-1.67709,0.037735) (-1.64134,0.039258) (-1.60559,0.040869) (-1.56983,0.042575) (-1.53408,0.044382) (-1.49832,0.046297) (-1.46257,0.048330) (-1.42682,0.050488) (-1.39106,0.052782) (-1.35531,0.055221) (-1.31955,0.057817) (-1.28380,0.060583) (-1.24804,0.063531) (-1.21229,0.066676) (-1.17654,0.070034) (-1.14078,0.073622) (-1.10503,0.077458) (-1.06927,0.081564) (-1.03352,0.085960) (-0.99777,0.090672) (-0.96201,0.095725) (-0.92626,0.101147) (-0.89050,0.106971) (-0.85475,0.113229) (-0.81899,0.119958) (-0.78324,0.127198) (-0.74749,0.134992) (-0.71173,0.143388) (-0.67598,0.152435) (-0.64022,0.162189) (-0.60447,0.172709) (-0.56872,0.184060) (-0.53296,0.196312) (-0.49721,0.209538) (-0.46145,0.223821) (-0.42570,0.239246) (-0.38994,0.255906) (-0.35419,0.273900) (-0.31844,0.293335) (-0.28268,0.314321) (-0.24693,0.336979) (-0.21117,0.361434) (-0.17542,0.387818) (-0.13966,0.416269) (-0.10391,0.446932) (-0.06816,0.479956) (-0.03240,0.515493) (0.00335,0.553700) (0.03911,0.594733) (0.07486,0.638752) (0.11061,0.685910) (0.14637,0.736357) (0.18212,0.790236) (0.21788,0.847678) (0.25363,0.908798) (0.28939,0.973691) (0.32514,1.042429) (0.36089,1.115053) (0.39665,1.191567) (0.43240,1.271932) (0.46816,1.356059) (0.50391,1.443804) (0.53966,1.534957) (0.57542,1.629234) (0.61117,1.726277) (0.64693,1.825637) (0.68268,1.926775) (0.71844,2.029055) (0.75419,2.131740) (0.78994,2.233992) (0.82570,2.334871) (0.86145,2.433339) (0.89721,2.528269) (0.93296,2.618454) (0.96872,2.702622) (1.00447,2.779462) (1.04022,2.847639) (1.07598,2.905834) (1.11173,2.952776) (1.14749,2.987280) (1.18324,3.008296) (1.21899,3.014955) (1.25475,3.006620) (1.29050,2.982935) (1.32626,2.943878) (1.36201,2.889802) (1.39777,2.821479) (1.43352,2.740125) (1.46927,2.647416) (1.50503,2.545493) (1.54078,2.436936) (1.57654,2.324730) (1.61229,2.212195) (1.64804,2.102900) (1.68380,2.000547) (1.71955,1.908834) (1.75531,1.831298) (1.79106,1.771145) (1.82682,1.731064) (1.86257,1.713059) (1.89832,1.718273) (1.93408,1.746857) (1.96983,1.797854) (2.00559,1.869154) (2.04134,1.957496) (2.07709,2.058541) (2.11285,2.167031) (2.14860,2.277011) (2.18436,2.382138) (2.22011,2.476043) (2.25587,2.552741) (2.29162,2.607058) (2.32737,2.635058) (2.36313,2.634413) (2.39888,2.604699) (2.43464,2.547578) (2.47039,2.466829) (2.50615,2.368210) (2.54190,2.259137) (2.57765,2.148196) (2.61341,2.044493) (2.64916,1.956893) (2.68492,1.893216) (2.72067,1.859441) (2.75642,1.859019) (2.79218,1.892370) (2.82793,1.956642) (2.86369,2.045788) (2.89944,2.150984) (2.93520,2.261397) (2.97095,2.365235) (3.00670,2.451002) (3.04246,2.508830) (3.07821,2.531726) (3.11397,2.516577) (3.14972,2.464759) (3.18547,2.382227) (3.22123,2.279016) (3.25698,2.168165) (3.29274,2.064152) (3.32849,1.981003) (3.36425,1.930323) (3.40000,1.919516)};\fill(0,.55)circle(1.5pt)node[left]{$1/4$};\draw[guide](-3,2.2)--(3.4,2.2);\node[right]at(3.4,2.2){$1$};}
无遮挡平面对应积分$-\infty<t<\infty$，值为$1+i$。因此图\ref{fig:o3edgea}使用的准确归一化为
\[
 \frac{U}{U_0}=\frac{(1/2+C)+i(1/2+S)}{1+i},\quad
 \frac I{I_0}=\frac{(1/2+C)^2+(1/2+S)^2}{2}.
\]
在$v=0$，分子为$(1+i)/2$，所以场为一半、强度为四分之一。$v\to-\infty$时$C,S\to-1/2$，深阴影强度趋零；$v\to+\infty$时均趋$1/2$，强度趋1。边界附近的空间尺度由$|v|\sim1$给出，即$|x|\sim\sqrt{\lambda L/2}$，所以屏越远，几何边缘被衍射展开的范围越大。正照、直边及均匀平面波是本归一化的条件。
\end{solution}

\begin{example}\label{ex:o3shiftproperty}
孔面场$U_0(\xi)$整体横移$\xi_0$。求夫琅禾费场的变化；若孔未移动而给孔面加线性相位$\exp(iq_0\xi)$，又怎样改变远场？
\Diagram{fig:o3shiftq}{孔的位置平移与孔内相位斜坡是两种不同操作}{\draw[axis](-1,0)--(5,0)node[right]{$\xi$};\draw[curve](0,0)rectangle(1.2,1);\draw[guide](2.5,0)rectangle(3.7,1);\draw[->](.6,1.5)--(3.1,1.5)node[midway,above]{$\xi_0$};}
\end{example}
\begin{solution}
取远场变量$q=k\sin\theta$，$F(q)=\int U_0(\xi)\ee^{-iq\xi}\dd\xi$。图\ref{fig:o3shiftq}平移孔为$U_0(\xi-\xi_0)$，令$s=\xi-\xi_0$，得$F_{\rm shift}(q)=\ee^{-iq\xi_0}F(q)$。单孔强度分布不变，只改变方向相关相位。
\Diagram{fig:o3shifta}{线性孔面相位使远场分布平移，而整体孔位置只加远场相位}{\draw[axis](-2.3,0)--(3.4,0)node[right]{$q$};\draw[curve,domain=-2:3,samples=70]plot(\x,{2*exp(-2*\x*\x)});\draw[guide,domain=-2:3,samples=70]plot(\x,{2*exp(-2*(\x-1.3)^2)});\node[above]at(1.3,2.1){$q_0$};}
若孔内乘$\ee^{iq_0\xi}$，积分变为$F(q-q_0)$，强度图样向$q_0$平移，如图\ref{fig:o3shifta}。斜入射照明的切向波矢或相位片的线性相位均按此式改变远场中心。孔整体平移在有另一参考孔相干叠加时会改变条纹，不能因单孔强度不变就删去其相位。
\end{solution}

\begin{example}\label{ex:o3twowidths}
两条中心相隔$b$的缝宽分别为$a_1,a_2$，单位面积照明场相同且无初始相差。求远场强度，说明可见度为何随角度变化。
\Diagram{fig:o3twowidthq}{不等缝宽带来不同角度包络和轴上振幅}{\draw[curve](0,-.6)rectangle(.6,1);\draw[curve](2,-.6)rectangle(3.2,1);\node[below]at(.3,-.6){$a_1$};\node[below]at(2.6,-.6){$a_2$};\draw[<->](.3,1.4)--(2.6,1.4)node[midway,above]{$b$};}
\end{example}
\begin{solution}
图\ref{fig:o3twowidthq}中每缝中心对称积分给实包络$F_j=a_j\operatorname{sinc}(\pi a_j\sin\theta/\lambda)$。合场$C[F_1+F_2\ee^{-ikb\sin\theta}]$，故$I=|C|^2[F_1^2+F_2^2+2F_1F_2\cos(kb\sin\theta)]$。
\Diagram{fig:o3twowidtha}{取$a_2=2a_1$时的有符号场包络，$\beta=\pi a_1\sin\theta/\lambda$}{
\draw[axis](-3.2,0)--(3.4,0)node[right]{$\beta$};
\draw[axis](0,-.65)--(0,2.7)node[above]{$F_j/a_1$};
\draw[curve,domain=-6.7:6.7,samples=181]plot({\x/2.3},{abs(\x)<.001 ? 1 : sin(deg(\x))/\x});
\draw[guide,domain=-6.7:6.7,samples=181]plot({\x/2.3},{abs(\x)<.001 ? 2 : sin(deg(2*\x))/\x});
\node[left]at(0,1){$1$};\node[left]at(0,2){$2$};
\node[below]at(1.366,0){$\pi$};\node[below]at(2.732,0){$2\pi$};
\draw[curve](1.6,2.2)--(2.2,2.2)node[right]{$F_1$};
\draw[guide](1.6,1.65)--(2.2,1.65)node[right]{$F_2$};}
在包络随一个纹距缓变的局部区域，$\mathcal V=2|F_1F_2|/(F_1^2+F_2^2)$。图\ref{fig:o3twowidtha}中两包络的零点不同：某一包络过零而另一不为零时，可见度降至0；过零后的$F_j$反号又使干涉纹相位翻转。若两者同时为零，该方向无光，不定义可见度。宽度比只给出轴上振幅比，离轴计算须保留各自的sinc因子。
\end{solution}

\section{连续孔径、二维光栅与采样}
\begin{example}\label{ex:o3circular}
半径$a$的均匀圆孔正入射。将远场场写成一维贝塞尔积分并求归一化形式，说明第一暗环如何决定角分辨尺度。
\Diagram{fig:o3circleq}{圆孔的径向对称性使远场只依赖离轴角}{\draw[curve](0,0)circle(1.6);\draw[axis](0,0)--(2.6,0)node[right]{$\xi$};\draw[axis](0,0)--(0,2.3)node[above]{$\eta$};\draw[->](0,0)--(1.1,1.1)node[right]{$r$};}
\end{example}
\begin{solution}
图\ref{fig:o3circleq}中，取观察方向的横向投影为极角零方向，$q=k\sin\theta$。孔面点的远场相位为$-qr\cos\varphi$，故
\[
 U=C\int_0^a r\dd r\int_0^{2\pi}\ee^{-iqr\cos\varphi}\dd\varphi
 =2\pi C\int_0^a rJ_0(qr)\dd r.
\]
这里用角积分定义$J_0(z)=(2\pi)^{-1}\int_0^{2\pi}\ee^{-iz\cos\varphi}\dd\varphi$。为求径向原函数，将指数展开；奇次余弦积分为零，偶次积分记为$C_{2j}=\int_0^{2\pi}\cos^{2j}\varphi\dd\varphi$。取$u=\cos^{2j-1}\varphi$、$\dd v=\cos\varphi\dd\varphi$，分部积分为
\begin{align*}
 C_{2j}&=[\cos^{2j-1}\varphi\sin\varphi]_0^{2\pi}
 +(2j-1)\int_0^{2\pi}\cos^{2j-2}\varphi\sin^2\varphi\dd\varphi\\
 &=(2j-1)(C_{2j-2}-C_{2j}),\qquad j\geq1.
\end{align*}
移项得$C_{2j}=(2j-1)C_{2j-2}/(2j)$，从$C_0=2\pi$递推，得$C_{2j}=2\pi(2j)!/[2^{2j}(j!)^2]$。于是
\[
 J_0(z)=\sum_{j=0}^{\infty}\frac{(-1)^j(z/2)^{2j}}{(j!)^2},
 \qquad
 J_1(z)=\sum_{j=0}^{\infty}
 \frac{(-1)^j(z/2)^{2j+1}}{j!(j+1)!}.
\]
第二式定义此处所需的一阶贝塞尔函数。逐项求导时，$zJ_1(z)$的第$j$项乘出$2j+2$，约去$(j+1)!$，便得到$\dd[zJ_1(z)]/\dd z=zJ_0(z)$。令$z=qr$、$r\dd r=z\dd z/q^2$，上下限由$0,a$变为$0,qa$，得到
\[
 U=\frac{2\pi C}{q^2}[zJ_1(z)]_0^{qa}
 =\frac{2\pi Ca}{q}J_1(qa).
\]
\Diagram{fig:o3circlea}{圆孔归一化远场强度；首零由贝塞尔函数确定}{
\draw[axis](0,0)--(5.1,0)node[right]{$qa$};
\draw[axis](0,0)--(0,2.9)node[above]{$I/I(0)$};
\draw[curve]plot coordinates{(0.000000,2.4500000) (0.015000,2.4494488) (0.030000,2.4477958) (0.045000,2.4450429) (0.060000,2.4411932) (0.075000,2.4362510) (0.090000,2.4302219) (0.105000,2.4231125) (0.120000,2.4149310) (0.135000,2.4056864) (0.150000,2.3953891) (0.165000,2.3840506) (0.180000,2.3716836) (0.195000,2.3583018) (0.210000,2.3439200) (0.225000,2.3285544) (0.240000,2.3122218) (0.255000,2.2949402) (0.270000,2.2767288) (0.285000,2.2576076) (0.300000,2.2375974) (0.315000,2.2167202) (0.330000,2.1949986) (0.345000,2.1724563) (0.360000,2.1491175) (0.375000,2.1250075) (0.390000,2.1001521) (0.405000,2.0745777) (0.420000,2.0483117) (0.435000,2.0213819) (0.450000,1.9938165) (0.465000,1.9656444) (0.480000,1.9368951) (0.495000,1.9075983) (0.510000,1.8777842) (0.525000,1.8474834) (0.540000,1.8167266) (0.555000,1.7855450) (0.570000,1.7539699) (0.585000,1.7220326) (0.600000,1.6897649) (0.615000,1.6571983) (0.630000,1.6243646) (0.645000,1.5912954) (0.660000,1.5580223) (0.675000,1.5245769) (0.690000,1.4909905) (0.705000,1.4572943) (0.720000,1.4235191) (0.735000,1.3896958) (0.750000,1.3558546) (0.765000,1.3220255) (0.780000,1.2882382) (0.795000,1.2545218) (0.810000,1.2209050) (0.825000,1.1874160) (0.840000,1.1540825) (0.855000,1.1209315) (0.870000,1.0879896) (0.885000,1.0552825) (0.900000,1.0228356) (0.915000,0.9906731) (0.930000,0.9588190) (0.945000,0.9272962) (0.960000,0.8961270) (0.975000,0.8653328) (0.990000,0.8349343) (1.005000,0.8049512) (1.020000,0.7754025) (1.035000,0.7463063) (1.050000,0.7176799) (1.065000,0.6895395) (1.080000,0.6619004) (1.095000,0.6347773) (1.110000,0.6081837) (1.125000,0.5821322) (1.140000,0.5566346) (1.155000,0.5317015) (1.170000,0.5073428) (1.185000,0.4835675) (1.200000,0.4603834) (1.215000,0.4377977) (1.230000,0.4158163) (1.245000,0.3944445) (1.260000,0.3736865) (1.275000,0.3535458) (1.290000,0.3340247) (1.305000,0.3151248) (1.320000,0.2968469) (1.335000,0.2791909) (1.350000,0.2621558) (1.365000,0.2457397) (1.380000,0.2299401) (1.395000,0.2147537) (1.410000,0.2001762) (1.425000,0.1862030) (1.440000,0.1728282) (1.455000,0.1600458) (1.470000,0.1478488) (1.485000,0.1362295) (1.500000,0.1251797) (1.515000,0.1146908) (1.530000,0.1047533) (1.545000,0.0953573) (1.560000,0.0864925) (1.575000,0.0781479) (1.590000,0.0703123) (1.605000,0.0629739) (1.620000,0.0561205) (1.635000,0.0497397) (1.650000,0.0438187) (1.665000,0.0383442) (1.680000,0.0333029) (1.695000,0.0286811) (1.710000,0.0244650) (1.725000,0.0206406) (1.740000,0.0171935) (1.755000,0.0141095) (1.770000,0.0113742) (1.785000,0.0089731) (1.800000,0.0068915) (1.815000,0.0051150) (1.830000,0.0036291) (1.845000,0.0024191) (1.860000,0.0014707) (1.875000,0.0007695) (1.890000,0.0003013) (1.905000,0.0000519) (1.920000,0.0000074) (1.935000,0.0001540) (1.950000,0.0004782) (1.965000,0.0009666) (1.980000,0.0016061) (1.995000,0.0023838) (2.010000,0.0032871) (2.025000,0.0043038) (2.040000,0.0054219) (2.055000,0.0066296) (2.070000,0.0079157) (2.085000,0.0092692) (2.100000,0.0106794) (2.115000,0.0121362) (2.130000,0.0136295) (2.145000,0.0151499) (2.160000,0.0166883) (2.175000,0.0182360) (2.190000,0.0197847) (2.205000,0.0213265) (2.220000,0.0228539) (2.235000,0.0243598) (2.250000,0.0258376) (2.265000,0.0272810) (2.280000,0.0286843) (2.295000,0.0300419) (2.310000,0.0313489) (2.325000,0.0326008) (2.340000,0.0337932) (2.355000,0.0349225) (2.370000,0.0359853) (2.385000,0.0369784) (2.400000,0.0378994) (2.415000,0.0387459) (2.430000,0.0395160) (2.445000,0.0402083) (2.460000,0.0408215) (2.475000,0.0413548) (2.490000,0.0418077) (2.505000,0.0421799) (2.520000,0.0424715) (2.535000,0.0426830) (2.550000,0.0428150) (2.565000,0.0428684) (2.580000,0.0428444) (2.595000,0.0427444) (2.610000,0.0425702) (2.625000,0.0423234) (2.640000,0.0420063) (2.655000,0.0416211) (2.670000,0.0411701) (2.685000,0.0406560) (2.700000,0.0400815) (2.715000,0.0394495) (2.730000,0.0387630) (2.745000,0.0380250) (2.760000,0.0372387) (2.775000,0.0364073) (2.790000,0.0355343) (2.805000,0.0346229) (2.820000,0.0336766) (2.835000,0.0326988) (2.850000,0.0316929) (2.865000,0.0306625) (2.880000,0.0296109) (2.895000,0.0285416) (2.910000,0.0274579) (2.925000,0.0263634) (2.940000,0.0252611) (2.955000,0.0241545) (2.970000,0.0230467) (2.985000,0.0219407) (3.000000,0.0208397) (3.015000,0.0197465) (3.030000,0.0186640) (3.045000,0.0175949) (3.060000,0.0165419) (3.075000,0.0155073) (3.090000,0.0144937) (3.105000,0.0135032) (3.120000,0.0125381) (3.135000,0.0116002) (3.150000,0.0106914) (3.165000,0.0098135) (3.180000,0.0089680) (3.195000,0.0081564) (3.210000,0.0073800) (3.225000,0.0066398) (3.240000,0.0059370) (3.255000,0.0052724) (3.270000,0.0046467) (3.285000,0.0040605) (3.300000,0.0035142) (3.315000,0.0030081) (3.330000,0.0025424) (3.345000,0.0021172) (3.360000,0.0017324) (3.375000,0.0013877) (3.390000,0.0010828) (3.405000,0.0008173) (3.420000,0.0005906) (3.435000,0.0004020) (3.450000,0.0002507) (3.465000,0.0001360) (3.480000,0.0000567) (3.495000,0.0000119) (3.510000,0.0000003) (3.525000,0.0000209) (3.540000,0.0000723) (3.555000,0.0001531) (3.570000,0.0002620) (3.585000,0.0003974) (3.600000,0.0005580) (3.615000,0.0007420) (3.630000,0.0009479) (3.645000,0.0011742) (3.660000,0.0014191) (3.675000,0.0016810) (3.690000,0.0019582) (3.705000,0.0022491) (3.720000,0.0025519) (3.735000,0.0028651) (3.750000,0.0031869) (3.765000,0.0035157) (3.780000,0.0038498) (3.795000,0.0041877) (3.810000,0.0045278) (3.825000,0.0048685) (3.840000,0.0052084) (3.855000,0.0055459) (3.870000,0.0058798) (3.885000,0.0062085) (3.900000,0.0065309) (3.915000,0.0068455) (3.930000,0.0071514) (3.945000,0.0074473) (3.960000,0.0077321) (3.975000,0.0080048) (3.990000,0.0082646) (4.005000,0.0085106) (4.020000,0.0087419) (4.035000,0.0089578) (4.050000,0.0091577) (4.065000,0.0093410) (4.080000,0.0095071) (4.095000,0.0096558) (4.110000,0.0097865) (4.125000,0.0098991) (4.140000,0.0099933) (4.155000,0.0100689) (4.170000,0.0101259) (4.185000,0.0101642) (4.200000,0.0101841) (4.215000,0.0101854) (4.230000,0.0101686) (4.245000,0.0101337) (4.260000,0.0100812) (4.275000,0.0100114) (4.290000,0.0099246) (4.305000,0.0098215) (4.320000,0.0097025) (4.335000,0.0095682) (4.350000,0.0094191) (4.365000,0.0092561) (4.380000,0.0090796) (4.395000,0.0088906) (4.410000,0.0086897) (4.425000,0.0084777) (4.440000,0.0082555) (4.455000,0.0080238) (4.470000,0.0077836) (4.485000,0.0075356) (4.500000,0.0072808)};
\node[left]at(0,2.45){$1$};
\draw[guide](1.915853,0)--(1.915853,.55);
\node[below]at(1.915853,0){$3.8317$};
\node[right]at(2.15,.45){首暗环};
}
轴上直接积分为$U(0)=C\pi a^2$，因此图\ref{fig:o3circlea}的归一化场和强度为
\[
 \frac{U}{U(0)}=\frac{2J_1(qa)}{qa},\qquad
 \frac I{I(0)}=\left[\frac{2J_1(qa)}{qa}\right]^2.
\]
$J_1(z)=z/2+O(z^3)$使轴上极限为1，不存在奇点。首个非零根$qa=3.8317$给
\[
 \sin\theta_0=\frac{3.8317}{2\pi}\frac\lambda a
 =1.2197\frac\lambda D,\quad D=2a.
\]
在小角度下这也是瑞利角间隔。两个互不相干点源的像应把两套强度相加；若相干，交叉项依相位而变，不能直接引用同一个双点强度判据。系数1.22由均匀圆孔的贝塞尔首零决定，矩形孔和高斯孔径需重新计算。
\end{solution}

\begin{example}\label{ex:o3gaussian}
孔面场振幅为$U(\xi)=U_0\ee^{-\xi^2/w^2}$，孔径足够大不截断高斯尾部。求远场强度及其$1/\ee^2$角半宽，与均匀硬边孔径比较。
\Diagram{fig:o3gaussq}{高斯振幅没有硬边界}{\draw[axis](-2.5,0)--(2.8,0)node[right]{$\xi/w$};\draw[axis](0,0)--(0,2.7)node[above]{$U/U_0$};\draw[curve,domain=-2.3:2.3,samples=80]plot(\x,{2.2*exp(-\x*\x)});\node[left]at(0,2.2){$1$};\draw(1,0)--(1,-.08)node[below]{$1$};}
\end{example}
\begin{solution}
图\ref{fig:o3gaussq}的远场积分为$F(q)=U_0\int_{-\infty}^{\infty}\ee^{-\xi^2/w^2-iq\xi}\dd\xi$。它正是第一讲已用分部积分及初值方程证明的式\eqref{eq:o1gaussintegral}。代入该式的$a=w^{-2}>0$、$b=-iq$，有$\sqrt{\pi/a}=w\sqrt\pi$、$b^2/(4a)=-q^2w^2/4$，因此
\[
 F(q)=U_0w\sqrt\pi\ee^{-q^2w^2/4}.
\]
轴上$q=0$时恢复孔面场的积分$U_0w\sqrt\pi$；方向相关的衰减完全由指数中的$q^2w^2$决定。
\Diagram{fig:o3gaussa}{高斯远场无旁瓣，角宽与孔面场宽成反比}{\draw[axis](-2.6,0)--(2.8,0)node[right]{$\theta/\theta_d$};\draw[axis](0,0)--(0,2.7)node[above]{$I/I(0)$};\draw[curve,domain=-2.4:2.4,samples=80]plot(\x,{2.3*exp(-2*\x*\x)});\draw[guide](-2.4,.3113)--(2.4,.3113);\node[left]at(0,.3113){$\ee^{-2}$};\node[left]at(0,2.3){$1$};\draw[guide](1,0)--(1,.3113);\draw(1,0)--(1,-.08)node[below]{$1$};}
图\ref{fig:o3gaussa}中的强度比为$\exp(-q^2w^2/2)$。取$I/I(0)=\ee^{-2}$，则$|q|=2/w$；近轴$q\simeq k\theta$，所以$1/\ee^2$角半宽为
\[
 \theta_d=\frac2{kw}=\frac\lambda{\pi w}.
\]
题中的$w$是场下降至$1/\ee$、强度下降至$1/\ee^2$的孔面半宽。若给的是强度高斯式$\exp(-\xi^2/w_I^2)$，须先取场的平方根，不能沿用相同的$w$。平滑高斯孔面对应无有限零点、无旁瓣的高斯远场，宽度须按约定强度比定义。$1/\ee^2$半宽与硬孔首零半宽采用不同的宽度标准。
\end{solution}

\begin{example}\label{ex:o3twoDgrating}
二维矩形周期阵列沿$x,y$周期为$b_x,b_y$，分别有$N_x,N_y$个相同小孔，正常入射。求主极大的方向条件及两方向首零宽度。
\Diagram{fig:o3twoDq}{二维有限孔阵列，图只显示部分周期}{\foreach\x in{0,1,2,3}{\foreach\y in{0,1,2}{\draw[curve](\x,\y)circle(.13);}}\draw[axis](-.5,-.5)--(4,-.5)node[right]{$x$};\draw[axis](-.5,-.5)--(-.5,3)node[above]{$y$};}
\end{example}
\begin{solution}
令$k_0=2\pi/\lambda$。图\ref{fig:o3twoDq}的平移相位可分离为
\[
 \sum_{j=0}^{N_x-1}\ee^{-ik_0b_xl_xj}
 \sum_{\ell=0}^{N_y-1}\ee^{-ik_0b_yl_y\ell}.
\]
两方向同时相长要求$b_xl_x=m\lambda$、$b_yl_y=n\lambda$。
\Diagram{fig:o3twoDa}{允许主峰位于方向余弦平面内的二维格点}{\draw[axis](-2,0)--(2.4,0)node[right]{$l_x$};\draw[axis](0,-1.8)--(0,2.2)node[above]{$l_y$};\draw[guide](0,0)circle(1.65);\foreach\x in{-1,0,1}{\foreach\y in{-1,0,1}{\fill(\x,\y)circle(1.5pt);}}}
还须$l_x^2+l_y^2\leq1$，图\ref{fig:o3twoDa}圆外格点不对应传播波。固定另一方向，首零宽度为$\Delta l_x=\lambda/(N_xb_x)$、$\Delta l_y=\lambda/(N_yb_y)$。单孔包络再乘到阵列因子上，某格点可能因孔形而缺级。二维条件不能只用总孔数$N_xN_y$替代一个方向的有效周期数。
\end{solution}

\begin{example}\label{ex:o3unequalresolution}
两条互不相干谱线的单线分布为偶函数$F$，强度比为$1:r$，峰中心为$\pm a$。解释为何等强时中点曲率判据不能直接用于$r\ne1$，写出真正峰合并边界的方程。
\Diagram{fig:o3unequalq}{弱谱线可能被强谱线的旁瓣或翼部覆盖}{\draw[axis](-3,0)--(3.3,0)node[right]{$x$};\draw[curve,domain=-2.8:2.8,samples=90]plot(\x,{2.3*exp(-2*(\x+1)^2)});\draw[guide,domain=-2.8:2.8,samples=90]plot(\x,{.7*exp(-2*(\x-1)^2)});}
\end{example}
\begin{solution}
图\ref{fig:o3unequalq}总强度$G(x)=F(x+a)+rF(x-a)$。当$r\ne1$时$G$不再关于0对称，$G'(0)=(1-r)F'(a)$一般不为零。因此即使$G''(0)=0$，中点也不是极值，不能用它判定两峰是否分开。
\Diagram{fig:o3unequala}{不等强高斯线形在临界间距处的合强度，$r=0.3$}{\draw[axis](-3.6,0)--(3.7,0)node[right]{$x$};\draw[axis](0,0)--(0,2.9)node[above]{$G$};\draw[curve]plot coordinates {(-3.40000,0.008830) (-3.35976,0.010673) (-3.31953,0.012859) (-3.27929,0.015443) (-3.23905,0.018487) (-3.19882,0.022059) (-3.15858,0.026235) (-3.11834,0.031102) (-3.07811,0.036752) (-3.03787,0.043289) (-2.99763,0.050823) (-2.95740,0.059476) (-2.91716,0.069376) (-2.87692,0.080663) (-2.83669,0.093484) (-2.79645,0.107991) (-2.75621,0.124347) (-2.71598,0.142717) (-2.67574,0.163271) (-2.63550,0.186182) (-2.59527,0.211622) (-2.55503,0.239760) (-2.51479,0.270761) (-2.47456,0.304783) (-2.43432,0.341970) (-2.39408,0.382454) (-2.35385,0.426349) (-2.31361,0.473745) (-2.27337,0.524709) (-2.23314,0.579277) (-2.19290,0.637453) (-2.15266,0.699204) (-2.11243,0.764459) (-2.07219,0.833103) (-2.03195,0.904978) (-1.99172,0.979877) (-1.95148,1.057549) (-1.91124,1.137690) (-1.87101,1.219952) (-1.83077,1.303937) (-1.79053,1.389206) (-1.75030,1.475273) (-1.71006,1.561619) (-1.66982,1.647686) (-1.62959,1.732891) (-1.58935,1.816630) (-1.54911,1.898282) (-1.50888,1.977220) (-1.46864,2.052818) (-1.42840,2.124458) (-1.38817,2.191542) (-1.34793,2.253499) (-1.30769,2.309791) (-1.26746,2.359927) (-1.22722,2.403466) (-1.18698,2.440025) (-1.14675,2.469287) (-1.10651,2.491006) (-1.06627,2.505012) (-1.02604,2.511211) (-0.98580,2.509591) (-0.94556,2.500220) (-0.90533,2.483249) (-0.86509,2.458905) (-0.82485,2.427492) (-0.78462,2.389386) (-0.74438,2.345025) (-0.70414,2.294912) (-0.66391,2.239598) (-0.62367,2.179680) (-0.58343,2.115789) (-0.54320,2.048583) (-0.50296,1.978736) (-0.46272,1.906932) (-0.42249,1.833851) (-0.38225,1.760163) (-0.34201,1.686518) (-0.30178,1.613541) (-0.26154,1.541821) (-0.22130,1.471903) (-0.18107,1.404289) (-0.14083,1.339426) (-0.10059,1.277707) (-0.06036,1.219464) (-0.02012,1.164969) (0.02012,1.114433) (0.06036,1.068004) (0.10059,1.025767) (0.14083,0.987750) (0.18107,0.953924) (0.22130,0.924203) (0.26154,0.898455) (0.30178,0.876497) (0.34201,0.858111) (0.38225,0.843037) (0.42249,0.830988) (0.46272,0.821651) (0.50296,0.814695) (0.54320,0.809771) (0.58343,0.806528) (0.62367,0.804608) (0.66391,0.803657) (0.70414,0.803329) (0.74438,0.803290) (0.78462,0.803223) (0.82485,0.802831) (0.86509,0.801841) (0.90533,0.800006) (0.94556,0.797108) (0.98580,0.792959) (1.02604,0.787405) (1.06627,0.780321) (1.10651,0.771619) (1.14675,0.761239) (1.18698,0.749154) (1.22722,0.735369) (1.26746,0.719914) (1.30769,0.702847) (1.34793,0.684251) (1.38817,0.664228) (1.42840,0.642900) (1.46864,0.620404) (1.50888,0.596890) (1.54911,0.572517) (1.58935,0.547451) (1.62959,0.521859) (1.66982,0.495912) (1.71006,0.469776) (1.75030,0.443616) (1.79053,0.417588) (1.83077,0.391839) (1.87101,0.366507) (1.91124,0.341720) (1.95148,0.317590) (1.99172,0.294219) (2.03195,0.271694) (2.07219,0.250087) (2.11243,0.229459) (2.15266,0.209856) (2.19290,0.191309) (2.23314,0.173839) (2.27337,0.157456) (2.31361,0.142157) (2.35385,0.127930) (2.39408,0.114756) (2.43432,0.102606) (2.47456,0.091446) (2.51479,0.081237) (2.55503,0.071934) (2.59527,0.063491) (2.63550,0.055858) (2.67574,0.048984) (2.71598,0.042817) (2.75621,0.037306) (2.79645,0.032398) (2.83669,0.028046) (2.87692,0.024200) (2.91716,0.020813) (2.95740,0.017843) (2.99763,0.015247) (3.03787,0.012987) (3.07811,0.011026) (3.11834,0.009331) (3.15858,0.007871) (3.19882,0.006618) (3.23905,0.005546) (3.27929,0.004633) (3.31953,0.003858) (3.35976,0.003202) (3.40000,0.002649)};\draw[guide](0.14049,0.803290)--(1.54049,0.803290);\fill(0.74049,0.803290)circle(1.5pt)node[above right]{$x_*$};}
应联立$F'(x_*+a)+rF'(x_*-a)=0$和$F''(x_*+a)+rF''(x_*-a)=0$，解未知$x_*,a$。例如$F(x)=\ee^{-x^2}$，两次求导为$F'=-2x\ee^{-x^2}$、$F''=(-2+4x^2)\ee^{-x^2}$。记$u=x_*+a,v=x_*-a$，第一式给$r\ee^{-v^2}=-(u/v)\ee^{-u^2}$；代入第二式并乘$v$，得到
\[
 v(4u^2-2)-u(4v^2-2)=2(u-v)(2uv+1)=0.
\]
非零分离$a>0$对应$uv=-1/2$，即$a^2-x_*^2=1/2$。再代回第一式，
\[
 r=\frac{a+x_*}{a-x_*}\ee^{-4ax_*}.
\]
对$r=0.3$取弱峰一侧$x_*>0$，数值求这两个标量方程，得$a=1.023881$、$x_*=0.740495$。图\ref{fig:o3unequala}按该解画出峰、谷合并处的水平切线。等强$r=1$时对称性给$x_*=0$，才退回中点曲率条件。
\end{solution}

\section{孔径自相关、成像传递与边缘加权}
\subsection{从点像积分得到强度传递函数}
在一个方向上讨论无像差、空间平移不变的近轴成像。设像面点响应的场为
\[
 h(x)=C\int P(\xi)\exp\!\left(-\frac{2\pi\mathrm i\xi x}{\lambda f}\right)\dd\xi,
\]
其中$P$为有效光瞳函数，$f$为此式所用的有效焦距尺度。对互不相干物点，归一化强度点响应是$g(x)=|h(x)|^2/\int|h|^2\dd x$。物面强度先换算到几何像坐标，再与$g$作卷积：
\[
 I(x)=\int I_{\rm geom}(x')g(x-x')\dd x'.
\]
这一步使用了不同物点间的交叉项平均为零，不能用来处理激光照明下的任意相干物体。

记未归一化强度的傅里叶变换为$N(\nu)$。将$|h|^2$展开，并对$x$积分，
\begin{align*}
 N(\nu)&=\int |h(x)|^2\ee^{-2\pi i\nu x}\dd x\\
 &=|C|^2\iint P(\xi)P^*(\eta)
 \left\{\int\ee^{-2\pi ix[(\xi-\eta)/(\lambda f)+\nu]}\dd x\right\}
 \dd\xi\dd\eta\\
 &=|C|^2\lambda f\int
 P(\xi)P^*(\xi+\lambda f\nu)\dd\xi.
\end{align*}
最后一步在傅里叶分布意义下使用$\int\ee^{-2\pi ixt}\dd x=\delta(t)$；括号内的狄拉克函数为$\lambda f\,\delta(\eta-\xi-\lambda f\nu)$，其尺度因子不能漏去。以$N(0)=|C|^2\lambda f\int|P|^2\dd\xi$归一化，得到
\[
 G(\nu)=\int g(x)\ee^{-2\pi\mathrm i\nu x}\dd x
 =\frac{\int P(\xi)P^*(\xi+\lambda f\nu)\dd\xi}
 {\int |P(\xi)|^2\dd\xi}.
\]
分母使$G(0)=1$。对于实对称光瞳，$G$为实偶函数；带像差的复光瞳还可能产生传递相位，不能只保留模。

宽$a$的均匀一维光瞳取$P=1$于$|\xi|<a/2$。分子是两段错开$\lambda f\nu$的区间重叠长度，故
\[
 G(\nu)=\begin{cases}1-|\nu|/\nu_c,&|\nu|\leq\nu_c,\\0,&|\nu|>\nu_c,
 \end{cases}\qquad \nu_c=\frac a{\lambda f}.
\]
上式给出不相干强度成像的截止频率。相干成像传递的是复场，传递函数与光瞳本身对应，截止为$a/(2\lambda f)$。比较这两个数值时须保留各自的照明条件和被传递的量；仅凭截止频率的比值，还不足以比较实际物体的可分辨程度。
\ExpansionDiagram{fig:dep-o3-overlap}{错位光瞳的重叠长度决定一维不相干传递函数}{
\fill[gray!20](2.5,.1)rectangle(5,1.5);\draw[very thick](1,.5)--(5,.5);\draw[very thick](2.5,1.2)--(6.5,1.2);\node[below]at(3,.1){重叠$a-|\lambda f\nu|$};\draw[<->](1,1.8)--(2.5,1.8)node[midway,above]{$\lambda f\nu$};}
图\ref{fig:dep-o3-overlap}同时说明了截止的来由：光瞳错位超过整个宽度后，重叠积分为零，物体相应的强度频率分量不再出现在像中。

\begin{example}条纹能够成像但对比度已经降低
某一维不相干成像系统的均匀有效光瞳宽$a=2.0\,\mathrm{mm}$，有效焦距$f=100\,\mathrm{mm}$，波长$\lambda=500\,\mathrm{nm}$。几何像的强度为$I_0[1+0.8\cos(2\pi\nu x)]$，其中$\nu=20\,\mathrm{mm}^{-1}$。求实际像的调制度，并求保持调制度不低于0.4所允许的最高频率。
\ExpansionDiagram{fig:dep-o3-modulation-q}{平移不变成像使各正弦强度分量分别乘传递系数}{
\draw[->](0,0)--(3.2,0)node[right]{$x$};\draw[->](0,0)--(0,2.1)node[above]{$I/I_0$};\draw[dotted](0,1)--(3,1);\draw[domain=0:3,samples=90]plot(\x,{1+.8*cos(360*\x)});\node[above]at(1.5,2.2){几何像};\draw[->](3.8,1)--(4.7,1);\draw[->](5,0)--(8.2,0)node[right]{$x$};\draw[->](5,0)--(5,2.1)node[above]{$I/I_0$};\draw[dotted](5,1)--(8,1);\draw[domain=5:8,samples=90]plot(\x,{1+.4*cos(360*\x)});\node[above]at(6.5,2.2){实际像};}
\end{example}
\begin{solution}
统一长度单位后，$\lambda=0.0005\,\mathrm{mm}$，所以$\nu_c=2/(0.0005\times100)=40\,\mathrm{mm}^{-1}$。常数分量乘$G(0)=1$，正弦分量乘$G(20)=1-20/40=1/2$，因此
\[
 I(x)=I_0[1+0.4\cos(2\pi\nu x)].
\]
图\ref{fig:dep-o3-modulation-q}中的周期不变，变化的是振幅相对平均强度的比例。由$I_{\max}=1.4I_0$、$I_{\min}=0.6I_0$，可见度也是0.4。
\ExpansionDiagram{fig:dep-o3-modulation-a}{调制度要求选出比截止频率更低的可用频段}{
\draw[->](0,0)--(5,0)node[right]{$\nu/\nu_c$};\draw[->](0,0)--(0,2.4)node[above]{$G$};\draw[thick](0,2)--(4,0);\draw[dashed](0,1)--(2,1)--(2,0);\node[left]at(0,2){1};\node[left]at(0,1){$1/2$};\node[below]at(2,0){$1/2$};\node[below]at(4,0){1};}
一般输出调制度为$0.8(1-\nu/\nu_c)$，要求不低于0.4给$\nu\leq\nu_c/2=20\,\mathrm{mm}^{-1}$。图\ref{fig:dep-o3-modulation-a}中截止前仍有一段低对比度区。因此“频率未超过截止”与“达到给定对比度”是两个不同条件。这里假设探测器能连续记录；有限像素还会另加像素孔径与采样限制。
\end{solution}

\subsection{连续加权为何能压低旁瓣}
硬边孔径在边缘突然截止，其傅里叶场具有较慢的尾部衰减。让透过振幅在边缘连续降至零，可以改变尾部阶数。判断代价时须固定照明条件，不能把峰值归一化后的两条曲线直接当作相同通光量。

\begin{example}三角振幅孔径与均匀孔径
一维孔宽均为$a$、入射场振幅同为$U_0$。第一孔的振幅透过率为1，第二孔为$t(\xi)=1-2|\xi|/a$于$|\xi|\leq a/2$，外侧为零。求第二孔的远场场、第一零点及总透过功率，与第一孔比较。
\ExpansionDiagram{fig:dep-o3-triangle-q}{相同外宽的硬边光瞳与三角振幅光瞳}{
\draw[->](-2,0)--(2,0)node[right]{$\xi$};\draw[dashed](-1.5,0)--(-1.5,1.5)--(1.5,1.5)--(1.5,0);\draw[thick](-1.5,0)--(0,1.5)--(1.5,0);\node[above]at(0,1.6){$t$};\node[below]at(-1.5,0){$-a/2$};\node[below]at(1.5,0){$a/2$};}
\end{example}
\begin{solution}
图\ref{fig:dep-o3-triangle-q}中三角透过率是两个宽$a/2$的单位矩形函数卷积后乘$2/a$。每个矩形的傅里叶变换为$(a/2)\operatorname{sinc}(qa/4)$，故卷积定理给
\[
 F_\triangle(q)=\frac{aU_0}{2}\operatorname{sinc}^2(qa/4),
 \quad \frac{I_\triangle(q)}{I_\triangle(0)}
 =\operatorname{sinc}^4(qa/4).
\]
也可直接积分核对。对$q\ne0$，分部积分取$u=1-2\xi/a$、$\dd v=\cos(q\xi)\dd\xi$，得
\begin{align*}
 F_\triangle(q)&=2U_0\int_0^{a/2}(1-2\xi/a)\cos(q\xi)\dd\xi\\
 &=2U_0\left[(1-2\xi/a)\frac{\sin(q\xi)}q\right]_0^{a/2}
 +\frac{4U_0}{aq}\int_0^{a/2}\sin(q\xi)\dd\xi\\
 &=\frac{4U_0}{aq}\left[-\frac{\cos(q\xi)}q\right]_0^{a/2}
 =\frac{4U_0}{aq^2}[1-\cos(qa/2)].
\end{align*}
边界项为零，最后用$1-\cos(qa/2)=2\sin^2(qa/4)$，重现卷积法结果。首零点为$|q|=4\pi/a$，均匀孔首零为$2\pi/a$，所以三角加权的中央瓣宽了一倍。
\ExpansionDiagram{fig:dep-o3-triangle-a}{归一化场的零点移动；实线为三角加权，虚线为均匀孔}{
\draw[->](-.2,0)--(7.2,0)node[right]{$qa/(2\pi)$};\draw[->](0,0)--(0,2.3)node[above]{归一化场};
\draw[thick,domain=.01:6.6,samples=180]plot(\x,{2*(sin(90*\x)/(pi*\x/2))^2});\draw[dashed,domain=.01:6.6,samples=180]plot(\x,{2*sin(180*\x)/(pi*\x)});\node[below=12pt]at(1,0){1};\node[below=12pt]at(2,0){2};}
图\ref{fig:dep-o3-triangle-a}画场而非强度，均匀孔场会反号，三角孔场为非负的平方。三角孔强度尾部按$q^{-4}$包络衰减，均匀孔按$q^{-2}$衰减，但主瓣展宽且透过功率降低。

功率需在孔面算$|t|^2$：
\[
 P_\triangle\propto2|U_0|^2\int_0^{a/2}(1-2\xi/a)^2\dd\xi
 =2|U_0|^2\left[\xi-\frac{2\xi^2}{a}+\frac{4\xi^3}{3a^2}\right]_0^{a/2}
 =\frac{a|U_0|^2}{3}.
\]
均匀孔为$a|U_0|^2$，故总功率比为$1/3$，轴上强度比则由场比$1/2$得$1/4$。这两个比值不同，原因是整个角分布的宽度和旁瓣都发生了改变。
\end{solution}

\section{孔面相位误差与相干峰的损失}
\subsection{误差的相关性先于误差的大小}
有$N$个等振幅、等间距单元时，远场场为$A(\delta)=a\sum_{j=0}^{N-1}\ee^{\mathrm i(j\delta+\epsilon_j)}$。相位误差全相同时，只有一个共同相位因子，强度完全不变；误差为$\epsilon_j=j\epsilon$时，分布整体移到$\delta=-\epsilon$；独立随机误差才会把部分窄峰能量转成背景。只报“均方相位误差相同”，不足以确定实际图样。

设各$\epsilon_j$均值为零，服从联合高斯分布，协方差为$C_{jk}=\langle\epsilon_j\epsilon_k\rangle$。场强平方含成对项，使用高斯特征函数有
\[
 \left\langle\ee^{\mathrm i(\epsilon_j-\epsilon_k)}\right\rangle
 =\exp[-\tfrac12(C_{jj}+C_{kk}-2C_{jk})].
\]
因而平均强度直接由每一对单元的相位差方差决定。共同误差使这个方差为零；远距离误差不相关时，其交叉项受到抑制。此式把完全相干阵列与部分相干孔径联系起来，也说明局部加工误差、整体倾斜和共同振动的影响不同。

若误差独立、方差同为$\sigma_\phi^2$，对角项始终为1，非对角项均为$\ee^{-\sigma_\phi^2}$，于是
\[
 \langle I(\delta)\rangle=|a|^2\left[
 \ee^{-\sigma_\phi^2}\left|\sum_j\ee^{\mathrm ij\delta}\right|^2
 +N(1-\ee^{-\sigma_\phi^2})\right].
\]
平均强度由等间距阵列的窄峰和角度均匀的背景组成，实际图样还要乘单元包络。在一个相位周期内积分，非对角指数项均为零，总积分仍是$2\pi N|a|^2$。这个纯相位模型没有吸收能量，峰值下降的同时，其他方向的强度增加。

\begin{example}相位误差对峰高和暗点的不同影响
取$N=100$，各单元独立高斯相位误差标准差为$\sigma_\phi=0.5$弧度。求平均主峰强度相对理想峰的比例、理想暗点处的平均背景，并求保持主峰至少为理想峰90\%所允许的最大$\sigma_\phi$。若误差来自正常入射反射面的独立高度误差，波长为$500\,\mathrm{nm}$，换算高度标准差。
\ExpansionDiagram{fig:dep-o3-phase-q}{单元位置不变，但各单元出射场的相位独立偏离}{
\draw[->](0,0)--(7,0)node[right]{单元};\foreach\x/\a in{.5/80,1.5/105,2.5/60,3.5/95,4.5/115,5.5/70}{\fill(\x,0)circle(1.2pt);\draw[->](\x,.3)--++(\a:.8);}\node[above]at(3,1.7){$\epsilon_j$改变相量方向};}
\end{example}
\begin{solution}
图\ref{fig:dep-o3-phase-q}中主峰的理想相位差为零，故
\[
 \frac{\langle I_{\max}\rangle}{N^2|a|^2}
 =\ee^{-\sigma_\phi^2}+
 \frac{1-\ee^{-\sigma_\phi^2}}N.
\]
代入$\ee^{-0.25}\simeq0.778801$，得到主峰比例0.781013。理想零点处的相干和为零，但背景为$N(1-\ee^{-0.25})|a|^2\simeq22.1199|a|^2$，不再严格消光。
\ExpansionDiagram{fig:dep-o3-phase-a}{独立相位误差降低平均主峰并填高理想零点；示意取10单元}{
\draw[->](0,0)--(6.8,0)node[right]{$\delta$};\draw[->](0,0)--(0,2.6)node[above]{强度};\draw[dashed,domain=.003:6.28,samples=400]plot(\x,{.022*(sin(5*\x r)/sin(.5*\x r))^2});\draw[thick,domain=.003:6.28,samples=400]plot(\x,{.022*(exp(-.25)*(sin(5*\x r)/sin(.5*\x r))^2+10*(1-exp(-.25)))});}
图\ref{fig:dep-o3-phase-a}用10单元显示曲线形状，本题数值计算使用100单元。要求主峰比例不低于0.9，等价于
\[
 \ee^{-\sigma_\phi^2}\geq\frac{0.9N-1}{N-1},\qquad
 \sigma_\phi\leq\sqrt{-\ln(89/99)}\simeq0.3263.
\]
正常入射反射面高度差造成往返程差$2\delta h$，故$\sigma_\phi=4\pi\sigma_h/\lambda$，允许$\sigma_h\lesssim13.0\,\mathrm{nm}$。若表面误差在各单元间强相关，不能继续用独立误差的背景公式；例如整面同步平移只改变共同相位，相干峰不因此降低。
\end{solution}

\subsection{连续孔径的平均相位与相位方差}
对均匀孔径，目标轴上场相对理想场为$\langle\ee^{\mathrm i\phi(\xi)}\rangle_P$，其中平均按孔径面积归一化。小相位误差展开至二阶：
\[
 \langle\ee^{\mathrm i\phi}\rangle_P
 \simeq1+\mathrm i\langle\phi\rangle_P-\frac12\langle\phi^2\rangle_P,
\]
所以轴上强度比
\[
 \mathcal S\simeq1-[\langle\phi^2\rangle_P-\langle\phi\rangle_P^2].
\]
共同相位从方差中抵消。这是围绕指定理想像点的二阶近似，须让去掉平均值后的相位足够小。若波前含线性倾斜，固定轴上点的强度会降低，但把观察点移到新峰中心后即可补偿；若含二次曲率，则需移动像面或改变聚焦条件。

\begin{example}一维光瞳的小离焦容许量
宽$a$的均匀一维孔径经焦距$f$的理想透镜聚焦。观察平面偏离后焦面$\delta$，$|\delta|\ll f$，忽略振幅传播因子的相对小变化。求轴上峰强相对该面理想相位补偿后峰强的二阶损失，并给出损失不超过5\%的近似离焦范围。
\ExpansionDiagram{fig:dep-o3-defocus-q}{同一孔径在偏离焦面的屏上留下二次相位}{
\draw[->](0,0)--(6,0)node[right]{$z$};\draw[<->,thick](.5,-1.4)--(.5,1.4);\draw(.5,1.2)--(4,0)--(5,-.34);\draw(.5,-1.2)--(4,0)--(5,.34);\draw[dashed](4,-1)--(4,1)node[above]{$f$};\draw[thick](5,-1)--(5,1)node[above]{$f+\delta$};}
\end{example}
\begin{solution}
图\ref{fig:dep-o3-defocus-q}中孔面到屏的近轴传播保留$kx^2/[2(f+\delta)]$，透镜原有相位为$-kx^2/(2f)$，故残余$\phi(x)=bx^2$，其中$b\simeq-k\delta/(2f^2)$。均匀区间$[-a/2,a/2]$的矩为
\[
 \langle x^2\rangle=\frac1a\int_{-a/2}^{a/2}x^2\dd x
 =\frac1a\left[\frac{x^3}{3}\right]_{-a/2}^{a/2}=\frac{a^2}{12},
\]
\[
 \langle x^4\rangle=\frac1a\int_{-a/2}^{a/2}x^4\dd x
 =\frac1a\left[\frac{x^5}{5}\right]_{-a/2}^{a/2}=\frac{a^4}{80},
 \quad \operatorname{Var}(x^2)=a^4\left(\frac1{80}-\frac1{144}\right)=\frac{a^4}{180}.
\]
于是
\[
 1-\mathcal S\simeq\frac{b^2a^4}{180}
 =\frac{k^2\delta^2a^4}{720f^4}.
\]
\ExpansionDiagram{fig:dep-o3-defocus-a}{去掉平均相位后的孔面相位仍呈弯曲，不能靠共同相移消除}{
\draw[->](-2.5,0)--(2.5,0)node[right]{$x$};\draw[->](0,-1.8)--(0,1.8)node[above]{$\phi-\langle\phi\rangle$};\draw[thick,domain=-2:2,samples=70]plot(\x,{-.5*((\x)^2-4/3)});\node[above]at(-2,0){$-a/2$};\node[above]at(2,0){$a/2$};}
图\ref{fig:dep-o3-defocus-a}表明影响峰强的是相位离散程度。要求上式不超过0.05，给
\[
 |\delta|\lesssim\frac{6f^2}{ka^2}=
 \frac{3\lambda f^2}{\pi a^2}
\]
（空气中$k=2\pi/\lambda$）。边界对应的均方根相位约$\sqrt{0.05}$弧度，仍属小误差估计；若要求更精确的阈值，应对完整菲涅耳积分求强度。几何弥散斑和轴上波动峰强是不同描述，不能仅把几何斑半径为零作为所有波前误差消失的条件。
\end{solution}

\section{习题与解答}
\begin{exercise}\label{ex:o3fresnelcircle}
均匀平面波正照半径$a$的圆孔，求距孔$L$的轴上菲涅耳场和强度。说明轴上为何可随孔径大小交替增强、减弱，不能一直用远场轴上场正比面积解释。
\Diagram{fig:o3fresnelcircleq}{孔内同心环带到轴上点的二次相位不同}{\draw[curve](0,0)ellipse(.5 and 1.6);\draw[guide](0,0)ellipse(.25 and .8);\fill(5,0)circle(1.5pt)node[right]{$P$};\draw[curve](0,1.6)--(5,0);\draw[curve](0,-1.6)--(5,0);\draw[guide](0,0)--(5,0);\node[above]at(2.5,.2){$L$};}
\end{exercise}
\begin{solution}
图\ref{fig:o3fresnelcircleq}轴上近轴传播积分为
\[
 U(P)=\frac{U_0\ee^{ikL}}{i\lambda L}
 2\pi\int_0^a r\ee^{ikr^2/(2L)}\dd r.
\]
令$q=kr^2/(2L)$，$r\dd r=(L/k)\dd q$，上下限变为$0,\phi_a$，其中$\phi_a=\pi a^2/(\lambda L)$。因此
\[
 \int_0^a r\ee^{ikr^2/(2L)}\dd r
 =\frac Lk\int_0^{\phi_a}\ee^{iq}\dd q
 =\frac Lk\left[\frac{\ee^{iq}}i\right]_0^{\phi_a}
 =\frac{L}{ik}(\ee^{i\phi_a}-1).
\]
代回传播前因子，并用$k\lambda=2\pi$，得$U(P)=U_0\ee^{ikL}[1-\ee^{i\phi_a}]$。所以相对无遮挡平面波强度
\[
 I(P)/I_0=4\sin^2(\phi_a/2).
\]
\Diagram{fig:o3fresnelcirclea}{轴上强度按边缘二次相位振荡}{\draw[axis](0,0)--(6.8,0)node[right]{$\phi_a$};\draw[axis](0,0)--(0,2.8)node[above]{$I/I_0$};\draw[curve,domain=0:6.283,samples=110]plot(\x,{2.3*sin((\x/2) r)^2});\node[left]at(0,2.3){$4$};\node[below]at(3.1416,0){$\pi$};\node[below]at(6.283,0){$2\pi$};}
图\ref{fig:o3fresnelcirclea}在$\phi_a\ll1$时$I/I_0\simeq\phi_a^2\propto a^4$，恢复远场轴上强度正比面积平方。较大孔径加入的环带可能与原场反相，使轴上减弱，能量分到轴外。积分仍要求近轴、孔径远大于波长；不能把$a\to\infty$的振荡直接当无限无遮挡平面波的极限，因为该极限不再一致满足所用截断条件。
\end{solution}

\begin{exercise}\label{ex:o3poisson}
把上一题圆孔换为同半径不透明圆盘，仍以平面波正照。用互补原理求圆盘后轴上菲涅耳场，解释几何阴影中心的亮点及其模型条件。
\Diagram{fig:o3poissonq}{圆盘遮挡几何直达光，边缘绕射仍到达轴上}{\draw[curve](0,0)ellipse(.45 and 1.5);\draw[very thick](0,-1.5)--(0,1.5);\draw[guide](0,1.5)--(5,0);\draw[guide](0,-1.5)--(5,0);\fill(5,0)circle(1.5pt)node[right]{$P$};\node[above]at(2.5,1.2){几何阴影内};}
\end{exercise}
\begin{solution}
图\ref{fig:o3poissonq}圆盘与圆孔互补，标量场线性给$U_{\rm disk}+U_{\rm hole}=U_{\rm free}$。上一题$U_{\rm hole}=U_{\rm free}(1-\ee^{i\phi_a})$，所以$U_{\rm disk}=U_{\rm free}\ee^{i\phi_a}$，轴上强度$I_{\rm disk}=I_0$。
\Diagram{fig:o3poissona}{轴上边缘环带的相位相同，形成阴影中央亮点}{\draw[curve](0,0)circle(1.5);\fill(0,0)circle(3pt);\foreach\a in{0,60,120,180,240,300}{\draw[guide]({1.5*cos(\a)},{1.5*sin(\a)})--(0,0);}\node[below]at(0,-1.9){观察屏中心示意};}
图\ref{fig:o3poissona}的中心亮点由圆周边缘到轴上等相位的衍射波叠加形成。实际有限尺寸光源、盘边误差、非轴对准和有限照明会降低或展宽亮点；理想强度结论限于轴上；阴影区其余位置仍可很暗，总透过功率也因圆盘遮挡而减小。
\end{solution}

\begin{exercise}\label{ex:o3babinet}
同一平面内两互补孔径的透过函数满足$t_1(\xi,\eta)+t_2(\xi,\eta)=1$。推导两者远场的关系，说明细丝与同宽狭缝何时具有相同衍射强度。
\Diagram{fig:o3babinetq}{细丝遮挡与同宽狭缝互补}{\begin{scope}[xshift=-2cm]\draw[curve](-1,-1.2)rectangle(1,1.2);\draw[very thick](0,-1.2)--(0,1.2);\node[below]at(0,-1.5){细丝};\end{scope}\begin{scope}[xshift=2cm]\draw[curve](-1,-1.2)--(-.15,-1.2)--(-.15,1.2)--(-1,1.2);\draw[curve](1,-1.2)--(.15,-1.2)--(.15,1.2)--(1,1.2);\node[below]at(0,-1.5){狭缝};\end{scope}}
\end{exercise}
\begin{solution}
图\ref{fig:o3babinetq}由衍射积分线性，$U_1+U_2=U_{\rm free}$。理想无限平面波的无遮挡远场只在前向方向有集中分量，因此离开前向时$U_2=-U_1$、$|U_2|^2=|U_1|^2$。
\Diagram{fig:o3babineta}{轴外两复场等大反相，强度相同}{\draw[axis](0,0)--(2.4,0)node[right]{$U_1$};\draw[axis](0,0)--(-2.4,0)node[left]{$U_2$};\node[above]at(0,1){轴外$U_{\rm free}=0$};}
所以细丝轴外暗纹可按同宽单缝公式测其宽度，如图\ref{fig:o3babineta}。但中央透过背景、有限照明包络和物体附近矢量边界场不相同，不能说两者整个屏幕总光强或总透过功率都相同。原理适用于相同入射波和标量孔径模型。
\end{solution}

\begin{exercise}\label{ex:o3phasehalves}
宽$a$狭缝的上、下半缝照明振幅相同，但相位相差$\pi$。以缝中心为原点，取场$+U_0$于$0<\xi<a/2$、$-U_0$于$-a/2<\xi<0$。求远场强度和零点。
\Diagram{fig:o3halvesq}{同一孔径两半具有反号复振幅}{\draw[curve](0,-1.5)rectangle(.7,1.5);\draw[curve](0,0)--(.7,0);\node[right]at(1,.8){$+U_0$};\node[right]at(1,-.8){$-U_0$};\node[left]at(0,0){$a$};}
\end{exercise}
\begin{solution}
图\ref{fig:o3halvesq}中上下半缝的场符号相反。以$q=k\sin\theta$，分别作两个定积分：
\begin{align*}
 F(q)/U_0&=\int_0^{a/2}\ee^{-iq\xi}\dd\xi
 -\int_{-a/2}^0\ee^{-iq\xi}\dd\xi\\
 &=\left[\frac{\ee^{-iq\xi}}{-iq}\right]_0^{a/2}
 -\left[\frac{\ee^{-iq\xi}}{-iq}\right]_{-a/2}^0\\
 &=\frac{1-\ee^{-iqa/2}}{iq}
 -\frac{\ee^{iqa/2}-1}{iq}\\
 &=\frac{2-2\cos(qa/2)}{iq}
 =-\frac{2i}{q}[1-\cos(qa/2)].
\end{align*}
故$I\propto4U_0^2[1-\cos(qa/2)]^2/q^2$。由于孔面场是实奇函数，远场为纯虚奇函数，也可由积分对称性预先判断。
\Diagram{fig:o3halvesa}{轴上因两半相消为暗，光能分到两侧瓣}{\draw[axis](-3.5,0)--(3.7,0)node[right]{$qa/2$};\draw[axis](0,0)--(0,2.7)node[above]{相对强度};\draw[curve,domain=-8:-.01,samples=180]plot({\x/2.5},{3.8*(1-cos(\x r))^2/(\x*\x)});\draw[curve,domain=.01:8,samples=180]plot({\x/2.5},{3.8*(1-cos(\x r))^2/(\x*\x)});\node[below]at(0,0){$0$};\node[below]at(2.51327,0){$2\pi$};\node[below]at(-2.51327,0){$-2\pi$};}
图\ref{fig:o3halvesa}的轴上必须取极限。展开$1-\cos(qa/2)=q^2a^2/8+O(q^4)$，得到$F\simeq-iU_0qa^2/4\to0$。对$q\ne0$，其他零点要求$qa/2=2m\pi$，即$a\sin\theta=2m\lambda$，$m\ne0$。

令$u=qa/4$，可写成$F=-iU_0a\sin^2u/u$。首个正旁瓣在$0<u<\pi$内，场模正比$\sin^2u/u$。用商法则，
\[
 \frac{\dd}{\dd u}\frac{\sin^2u}{u}
 =\frac{2u\sin u\cos u-\sin^2u}{u^2}.
\]
内点$\sin u\ne0$，驻值方程为$2u\cos u=\sin u$，即$\tan u=2u$；取该区间的非零根。相位片没有改变孔面强度，按帕塞瓦尔关系总远场功率保持不变，原中央能量转移到两侧。
\end{solution}

\begin{exercise}\label{ex:o3sinusoidalphase}
在近轴薄元件的标量模型中，无限宽相位光栅透过率为$t(x)=\exp[i\phi\cos(2\pi x/b)]$。设正入射，$\lambda/b\ll1$，忽略反射及各低级衍射束方向余弦的差别。定义第$m$级的傅里叶系数$c_m$，在$|\phi|\ll1$时求零级、正负一级振幅到所需阶，并验证这一模型中的功率守恒到二阶。
\Diagram{fig:o3sinephaseq}{正弦相位调制不改变孔面振幅}{\draw[axis](0,0)--(5,0)node[right]{$x$};\draw[axis](0,-1.4)--(0,1.7)node[above]{相位};\draw[curve,domain=0:4.5,samples=100]plot(\x,{cos(240*\x)});\node[below]at(2.7,-1.4){周期$b$};}
\end{exercise}
\begin{solution}
图\ref{fig:o3sinephaseq}的周期函数展开为$t=\sum_mc_m\ee^{i2\pi mx/b}$，系数为
\[
 c_m=\frac1b\int_0^bt(x)\ee^{-i2\pi mx/b}\dd x.
\]
令$u=2\pi x/b$，弱调制展开为$t\simeq1+i\phi\cos u-\phi^2\cos^2u/2$。
\Diagram{fig:o3sinephasea}{弱相位光栅把少量零级功率转入两侧一级}{\draw[axis](-2.5,0)--(2.7,0)node[right]{$m$};\draw[very thick](0,0)--(0,2.4);\draw[very thick](-1.5,0)--(-1.5,.6);\draw[very thick](1.5,0)--(1.5,.6);\node[below]at(-1.5,0){$-1$};\node[below]at(1.5,0){$1$};}
用$\cos u=(\ee^{iu}+\ee^{-iu})/2$和$\cos^2u=(1+\cos2u)/2$，逐项收集指数函数：
\[
 t=1-\frac{\phi^2}{4}
 +\frac{i\phi}{2}(\ee^{iu}+\ee^{-iu})
 -\frac{\phi^2}{8}(\ee^{2iu}+\ee^{-2iu})+O(\phi^3).
\]
因此$c_0=1-\phi^2/4+O(\phi^4)$，$c_{\pm1}=i\phi/2+O(\phi^3)$；
二级振幅虽为$O(\phi^2)$，其功率仅从四阶起算。在题设的近轴归一化下，图\ref{fig:o3sinephasea}中三束的功率份额分别为
\[
 \eta_0=|c_0|^2=1-\frac{\phi^2}{2}+O(\phi^4),\qquad
 \eta_{\pm1}=|c_{\pm1}|^2=\frac{\phi^2}{4}+O(\phi^4).
\]
相加得到$\eta_0+\eta_{+1}+\eta_{-1}=1+O(\phi^4)$，即守恒到二阶。
图中竖线只标示三束分配，不按小参数的真实比例绘制；弱调制时一级功率远小于零级。

传播角满足$\sin\theta_m=m\lambda/b$，题设保证所保留的一级为近轴传播波。
若衍射角较大，以电场幅值计量时法向能流还含方向余弦因子；若某级不传播，它也不能作为远场输出功率相加。
大角度情形的能流分配应由实际边界条件和各级传播方向确定。孔面恒模条件$|t|=1$给出傅里叶系数的平方和，却不足以单独确定这些法向能流。相位调制较大时，还须保留更多衍射级。
\end{solution}

\begin{exercise}\label{ex:o3camera}
直径$D=50$ mm、焦距$f=200$ mm的理想圆孔成像系统在550 nm工作。求后焦面点像首暗环半径与瑞利角间隔，并与宽度同为50 mm的长直缝比较。
\Diagram{fig:o3cameraq}{圆孔点像落在后焦面}{\OptAxis\draw[<->](1,-1.3)--(1,1.3);\draw[curve](5,-1.6)--(5,1.6);\draw[curve](1,1.2)--(5,0);\draw[curve](1,-1.2)--(5,0);\node[above]at(1,1.3){$D$};\node[above]at(5,1.6){焦面};}
\end{exercise}
\begin{solution}
图\ref{fig:o3cameraq}圆孔首暗角$\theta_0\simeq1.22\lambda/D=1.342\times10^{-5}$弧度，焦面半径$\rho_0=f\theta_0\simeq2.684\,\mu\mathrm m$。
\Diagram{fig:o3cameraa}{圆孔圆斑半径与长缝单方向宽度使用不同系数}{\draw[curve](-1.5,0)circle(.7);\draw[curve](1.5,-1)--(1.5,1);\draw[curve](2.1,-1)--(2.1,1);\node[below]at(-1.5,-1.3){圆孔};\node[below]at(1.8,-1.3){长缝};}
同宽长缝首暗角$\lambda/a=1.10\times10^{-5}$弧度，焦面半宽为$2.20\,\mu\mathrm m$。图\ref{fig:o3cameraa}比较的都是中心到首零，圆孔有1.22因子，长缝没有。两者通光面积、二维像形和旁瓣分布不同，不能仅以一个更小数值判定所有成像指标更好。
\end{solution}

\begin{exercise}\label{ex:o3missinggeneral}
光栅周期与缝宽之比$b/a=p/q$，其中$p,q$为互素正整数且$p>q$。求缺级的全部整数主极大级次。
\Diagram{fig:o3missinggeneralq}{缺级要求光栅主峰和单缝暗纹重合}{\draw[axis](-.3,0)--(5,0)node[right]{级次};\foreach\x in{0,1,2,3,4}{\draw[curve](\x,0)--(\x,1);}\draw[guide](0,1.8)..controls(2,2.5)and(3,.1)..(4.5,0);\node[right]at(4.5,1.6){包络零点};}
\end{exercise}
\begin{solution}
图\ref{fig:o3missinggeneralq}重合条件$b\sin\theta=m\lambda$、$a\sin\theta=s\lambda$，$s\ne0$。相除$m/s=p/q$，所以$q m=p s$。由$p,q$互素，$m=kp,s=kq$，$k$为非零整数。
\Diagram{fig:o3missinggenerala}{最先缺掉的正级次为$p$，随后为其整数倍}{\draw[axis](0,0)--(6.7,0)node[right]{$m$};\foreach\x in{1,2,3,4,5,6}{\draw[curve](\x,0)--(\x,.8);}\foreach\x in{3,6}{\draw[thick](\x-.2,.5)--(\x+.2,.9);\draw[thick](\x-.2,.9)--(\x+.2,.5);}\node[above]at(3,1.2){例$p=3$};}
故缺级为$\pm p,\pm2p,\ldots$，前提这些级次本来可传播。图\ref{fig:o3missinggenerala}显示$p=3$例。若$b/a$为无理数，理想精确模型没有完全重合的非零级次，但可能有非常接近的弱主峰，实验上是否可见还取决于包络和动态范围。
\end{solution}

\begin{exercise}\label{ex:o3exposure}
双缝中一缝被振动镜调制相位$\phi(t)=\phi_a\cos\Omega t$。给出长曝光干涉可见度的积分表达式，并在$\phi_a\ll1$时求二阶近似。
\Diagram{fig:o3exposureq}{一臂相位作周期调制}{\draw[->](0,.8)--(4.5,.8);\draw[->](0,-.3)--(4.5,-.3);\draw[very thick](2,-.8)--(2,.2);\draw[<->](1.7,-1.3)--(2.3,-1.3);\node[below]at(2,-1.5){振动元件};}
\end{exercise}
\begin{solution}
图\ref{fig:o3exposureq}中两束等强，固定屏点的瞬时强度为$2I_1[1+\cos(\delta_0+\phi_a\cos\Omega t)]$。长曝光覆盖许多周期，可将时间平均改成对$u=\Omega t$的周期平均，
\[
 G=\frac1{2\pi}\int_0^{2\pi}\ee^{i\phi_a\cos u}\dd u=J_0(\phi_a).
\]
虚部为零，因为$u$与$u+\pi$的正弦项相反；实部保留。平均强度为$2I_1[1+J_0(\phi_a)\cos\delta_0]$，可见度是$|J_0(\phi_a)|$。
\Diagram{fig:o3exposurea}{小幅调制的平均相量比单位长度缩短}{\draw[curve](0,0)circle(1.4);\draw[axis](0,0)--(1.4,0);\draw[guide](0,0)--(1.2,.72);\draw[guide](0,0)--(1.2,-.72);\draw[very thick,->](0,0)--(1.15,0);\node[below]at(0,-1.8){平均相量沿实轴};}
图\ref{fig:o3exposurea}表示相量绕静态相位来回摆动后的平均。使用例\ref{ex:o3circular}已从角积分推导的$J_0$级数，分别取$j=0,1,2$项，得到
\[
 J_0(\phi_a)=1-\frac{\phi_a^2}{4}+\frac{\phi_a^4}{64}+\cdots.
\]
到二阶$\mathcal V\simeq1-\phi_a^2/4$。若振幅增大使$J_0$越过零点，条纹先消失再出现，且亮暗相位翻转$\pi$；因此大振幅下不能把二阶式外推成负可见度。相位方差虽为$\phi_a^2/2$，其分布并非高斯，直接套高斯相位噪声公式只会在最低非平凡阶巧合一致。短曝光还应保留实际起始相位和曝光区间。
\end{solution}

\begin{exercise}\label{ex:o3gratingorderoverlap}
光栅中波长$\lambda_1$的$m$级与波长$\lambda_2$的$(m+1)$级方向重合。求两波长关系，说明增加$N$能否分开这种重级。
\Diagram{fig:o3overlapq}{不同波长、不同级次可以具有同一出射角}{\draw[curve](0,-1.2)--(0,1.2);\draw[axis](0,0)--(3.7,1.6)node[right]{共同方向};\draw[guide](0,0)--(4,0);\node[above]at(1.9,1.8){$m\lambda_1=(m+1)\lambda_2$};}
\end{exercise}
\begin{solution}
图\ref{fig:o3overlapq}同角要求$b\sin\theta=m\lambda_1=(m+1)\lambda_2$，所以$\lambda_2=m\lambda_1/(m+1)$。两个理想峰中心完全重合，即使$N$增大、各峰变窄，也不会产生角位置分离。
\Diagram{fig:o3overlapa}{同中心的两峰变窄后仍重叠，需要级次筛选或额外色散}{\draw[axis](-2.3,0)--(2.5,0)node[right]{$\theta-\theta_0$};\draw[curve,domain=-2:2,samples=80]plot(\x,{2.2*exp(-4*\x*\x)});\draw[guide,domain=-2:2,samples=80]plot(\x,{1.5*exp(-8*\x*\x)});}
图\ref{fig:o3overlapa}说明同级分辨率$mN$不能处理不同级完全重叠。可用滤光片限制输入谱段，或用棱镜等另一色散方向交叉分散。若两波长的单孔包络不同，强度可不同，但仅凭同一位置的强度不能唯一拆出两个未知谱线。
\end{solution}

\begin{exercise}\label{ex:o3illumination}
总周期数$N$的光栅只有中间$N/2$条被均匀照亮，其余无光。求实际分辨率和主峰宽度相对全部照明的变化，讨论保持总入射功率与保持每缝照明两种条件。
\Diagram{fig:o3illumq}{实体缝数与实际照明缝数不同}{\draw[curve](0,-.6)rectangle(5,.8);\foreach\x in{.25,.75,1.25,1.75,2.25,2.75,3.25,3.75,4.25,4.75}{\draw[curve](\x,-.6)--(\x,.8);}\draw[guide](1.25,-1)--(1.25,1.2);\draw[guide](3.75,-1)--(3.75,1.2);\node[above]at(2.5,1.2){照明区域};}
\end{exercise}
\begin{solution}
图\ref{fig:o3illumq}实际相干求和只有$N_{\rm eff}=N/2$项，首零角半宽变为原来2倍，瑞利分辨率$mN/2$。计算应采用实际受到相干照明的条数。
\Diagram{fig:o3illuma}{缩小照明宽度加宽主峰；纵轴高度取决于功率约束}{\draw[axis](-2.5,0)--(2.7,0)node[right]{角偏移};\draw[curve,domain=-2.3:2.3,samples=80]plot(\x,{2.3*exp(-4*\x*\x)});\draw[guide,domain=-2.3:2.3,samples=80]plot(\x,{1.15*exp(-\x*\x)});}
若每缝场不变，主峰强度按项数平方降为1/4，总透过功率降为1/2；若把原总功率集中到半数缝，每缝强度加倍，则主峰降为1/2，总功率不变。图\ref{fig:o3illuma}采用后一种归一化示意。角宽结论不受这两种整体振幅标定的影响。
\end{solution}

\begin{exercise}\label{ex:o3tiltblaze}
反射闪耀光栅的入射方向相对宏观外法线与出射方向分居两侧，角度大小为$i,\theta$。小面法线相对宏观法线倾角$\alpha$。写出小面镜面条件和光栅相长条件，说明为何必须先画方向再定正负。
\Diagram{fig:o3tiltblazeq}{入射与出射分居法线两侧，采用角度大小约定}{\draw[curve](-2.5,0)--(2.5,0);\draw[guide](0,0)--(0,2.7);\draw[->](-1.5,2.1)--(0,0);\draw[->](0,0)--(2,1.7);\draw[guide](0,0)--(.6,2.6);\node[above]at(.6,2.6){小面法线};\node[left]at(-.7,1.3){$i$};\node[right]at(1,1){$\theta$};}
\end{exercise}
\begin{solution}
按图\ref{fig:o3tiltblazeq}，从表面指向光源的方向角为$-i$，出射为$+\theta$。镜面法线平分二者，故$\theta-i=2\alpha$。入射实际波矢沿向下向右，切向分量为$+k\sin i$；出射切向为$+k\sin\theta$，所以相长条件可取$b(\sin\theta-\sin i)=m\lambda$。
\Diagram{fig:o3tiltblazea}{以实际波矢切向分量之差确定级次，法线平分反向入射与出射}{\draw[axis](-2.2,0)--(2.8,0)node[right]{切向};\draw[axis](0,0)--(1.4,2)node[right]{$\vct k_o$};\draw[axis](-1.6,2)--(0,0)node[midway,left]{$\vct k_i$};\draw[guide](0,0)--(.4,2.6);}
联立$\theta=i+2\alpha$，得$\lambda_B=b[\sin(i+2\alpha)-\sin i]/m$。正常入射恢复$b\sin2\alpha/m$。图\ref{fig:o3tiltblazea}的角度约定不同于有些把入射与出射都按波矢同侧有向角定义的写法；公式应从图中切向波矢推得，不能把两个正弦的加减号脱离图使用。
\end{solution}

\begin{exercise}\label{ex:o3farfield}
宽$a=0.2$ mm的缝受$\lambda=500$ nm平面波照射，屏距$L=1$ m。估计菲涅耳数，求首暗纹位置，并说明把宽度加倍对近似条件与主瓣宽度的影响。
\Diagram{fig:o3farq}{有限宽缝到远屏的几何尺度}{\draw[curve](0,-1.2)--(0,-.2);\draw[curve](0,.2)--(0,1.2);\draw[curve](5,-1.5)--(5,1.8);\draw[guide](0,0)--(5,0);\node[above]at(2.5,0){$L$};\node[left]at(0,0){$a$};}
\end{exercise}
\begin{solution}
用全宽定义估计数$F_a=a^2/(\lambda L)=.08$；若以半宽定义，数值为.02。图\ref{fig:o3farq}孔边相对中心的二次相位是$\pi a^2/(4\lambda L)=.0628$弧度，足够小，夫琅禾费近似合理。
\Diagram{fig:o3fara}{宽缝使主瓣缩窄，同时加大被省略的二次相位}{\draw[axis](-2.7,0)--(2.9,0)node[right]{$x$};\draw[curve](-2,0)..controls(-.7,2.8)and(.7,2.8)..(2,0);\draw[guide](-1,0)..controls(-.4,2.8)and(.4,2.8)..(1,0);\node[above]at(0,2.4){$a$与$2a$示意};}
首暗纹$\sin\theta=\lambda/a=.0025$，屏坐标$x=L\tan\theta\simeq2.5$ mm。宽度加倍后首暗纹约1.25 mm，但菲涅耳数增为4倍，二次相位.251弧度，远场近似精度下降。图\ref{fig:o3fara}只比较主瓣宽度；若孔面照明强度固定，轴上强度还会增加，不能把归一化曲线当绝对光强不变。
\end{solution}

\begin{exercise}\label{ex:o3energy}
均匀平面波照亮单缝，缝宽由$a$增到$2a$。为何轴上强度增为4倍，总透过功率却只增为2倍？用积分说明。
\Diagram{fig:o3energyq}{缝宽变化，孔面单位面积照明保持不变}{\draw[curve](0,-1)rectangle(.6,1);\draw[curve](3,-1)rectangle(4.2,1);\node[below]at(.3,-1){$a$};\node[below]at(3.6,-1){$2a$};\draw[->](1,0)--(2.6,0);}
\end{exercise}
\begin{solution}
图\ref{fig:o3energyq}中，孔面场强固定，轴上各点同相，所以$U(0)\propto a$、$I(0)\propto a^2$。设一维孔函数$g(\xi)=1$于$|\xi|<a/2$，其他为零，其傅里叶变换$F(q)=a\operatorname{sinc}(qa/2)$。帕塞瓦尔关系给
\[
 \int_{-\infty}^{\infty}|F(q)|^2\frac{\dd q}{2\pi}
 =\int_{-\infty}^{\infty}|g(\xi)|^2\dd\xi=a.
\]
令$u=qa/2$，得到$\int_{-\infty}^{\infty}\operatorname{sinc}^2u\dd u=\pi$，也明确了所用积分常数。
\Diagram{fig:o3energya}{同一照明下的准确单缝强度；宽度加倍使轴上峰高增为4倍}{
\draw[axis](-3.1,0)--(3.3,0)node[right]{$u$};
\draw[axis](0,0)--(0,2.65)node[above]{$I/I_a(0)$};
\draw[curve]plot coordinates{(-2.800000,0.0051091) (-2.786667,0.0047656) (-2.773333,0.0044207) (-2.760000,0.0040761) (-2.746667,0.0037336) (-2.733333,0.0033951) (-2.720000,0.0030624) (-2.706667,0.0027374) (-2.693333,0.0024218) (-2.680000,0.0021177) (-2.666667,0.0018269) (-2.653333,0.0015511) (-2.640000,0.0012923) (-2.626667,0.0010523) (-2.613333,0.0008327) (-2.600000,0.0006353) (-2.586667,0.0004617) (-2.573333,0.0003136) (-2.560000,0.0001923) (-2.546667,0.0000995) (-2.533333,0.0000363) (-2.520000,0.0000041) (-2.506667,0.0000040) (-2.493333,0.0000371) (-2.480000,0.0001042) (-2.466667,0.0002061) (-2.453333,0.0003436) (-2.440000,0.0005172) (-2.426667,0.0007273) (-2.413333,0.0009741) (-2.400000,0.0012578) (-2.386667,0.0015784) (-2.373333,0.0019356) (-2.360000,0.0023291) (-2.346667,0.0027584) (-2.333333,0.0032228) (-2.320000,0.0037216) (-2.306667,0.0042538) (-2.293333,0.0048182) (-2.280000,0.0054136) (-2.266667,0.0060385) (-2.253333,0.0066912) (-2.240000,0.0073702) (-2.226667,0.0080734) (-2.213333,0.0087988) (-2.200000,0.0095443) (-2.186667,0.0103076) (-2.173333,0.0110863) (-2.160000,0.0118778) (-2.146667,0.0126794) (-2.133333,0.0134885) (-2.120000,0.0143022) (-2.106667,0.0151175) (-2.093333,0.0159316) (-2.080000,0.0167413) (-2.066667,0.0175436) (-2.053333,0.0183353) (-2.040000,0.0191133) (-2.026667,0.0198744) (-2.013333,0.0206154) (-2.000000,0.0213332) (-1.986667,0.0220247) (-1.973333,0.0226867) (-1.960000,0.0233163) (-1.946667,0.0239104) (-1.933333,0.0244662) (-1.920000,0.0249809) (-1.906667,0.0254518) (-1.893333,0.0258764) (-1.880000,0.0262522) (-1.866667,0.0265770) (-1.853333,0.0268487) (-1.840000,0.0270654) (-1.826667,0.0272254) (-1.813333,0.0273271) (-1.800000,0.0273693) (-1.786667,0.0273508) (-1.773333,0.0272709) (-1.760000,0.0271290) (-1.746667,0.0269247) (-1.733333,0.0266581) (-1.720000,0.0263293) (-1.706667,0.0259389) (-1.693333,0.0254878) (-1.680000,0.0249770) (-1.666667,0.0244080) (-1.653333,0.0237825) (-1.640000,0.0231026) (-1.626667,0.0223707) (-1.613333,0.0215896) (-1.600000,0.0207622) (-1.586667,0.0198919) (-1.573333,0.0189824) (-1.560000,0.0180377) (-1.546667,0.0170621) (-1.533333,0.0160602) (-1.520000,0.0150370) (-1.506667,0.0139978) (-1.493333,0.0129480) (-1.480000,0.0118934) (-1.466667,0.0108403) (-1.453333,0.0097949) (-1.440000,0.0087637) (-1.426667,0.0077538) (-1.413333,0.0067721) (-1.400000,0.0058260) (-1.386667,0.0049228) (-1.373333,0.0040704) (-1.360000,0.0032764) (-1.346667,0.0025488) (-1.333333,0.0018957) (-1.320000,0.0013253) (-1.306667,0.0008458) (-1.293333,0.0004656) (-1.280000,0.0001930) (-1.266667,0.0000364) (-1.253333,0.0000040) (-1.240000,0.0001043) (-1.226667,0.0003456) (-1.213333,0.0007359) (-1.200000,0.0012834) (-1.186667,0.0019960) (-1.173333,0.0028816) (-1.160000,0.0039476) (-1.146667,0.0052016) (-1.133333,0.0066507) (-1.120000,0.0083018) (-1.106667,0.0101615) (-1.093333,0.0122363) (-1.080000,0.0145321) (-1.066667,0.0170546) (-1.053333,0.0198090) (-1.040000,0.0228003) (-1.026667,0.0260330) (-1.013333,0.0295110) (-1.000000,0.0332381) (-0.986667,0.0372172) (-0.973333,0.0414511) (-0.960000,0.0459419) (-0.946667,0.0506911) (-0.933333,0.0556999) (-0.920000,0.0609687) (-0.906667,0.0664974) (-0.893333,0.0722856) (-0.880000,0.0783319) (-0.866667,0.0846346) (-0.853333,0.0911913) (-0.840000,0.0979990) (-0.826667,0.1050541) (-0.813333,0.1123524) (-0.800000,0.1198892) (-0.786667,0.1276589) (-0.773333,0.1356556) (-0.760000,0.1438728) (-0.746667,0.1523031) (-0.733333,0.1609388) (-0.720000,0.1697716) (-0.706667,0.1787925) (-0.693333,0.1879921) (-0.680000,0.1973604) (-0.666667,0.2068868) (-0.653333,0.2165603) (-0.640000,0.2263693) (-0.626667,0.2363020) (-0.613333,0.2463458) (-0.600000,0.2564879) (-0.586667,0.2667151) (-0.573333,0.2770136) (-0.560000,0.2873696) (-0.546667,0.2977687) (-0.533333,0.3081964) (-0.520000,0.3186377) (-0.506667,0.3290776) (-0.493333,0.3395008) (-0.480000,0.3498918) (-0.466667,0.3602350) (-0.453333,0.3705147) (-0.440000,0.3807151) (-0.426667,0.3908205) (-0.413333,0.4008149) (-0.400000,0.4106826) (-0.386667,0.4204078) (-0.373333,0.4299750) (-0.360000,0.4393687) (-0.346667,0.4485734) (-0.333333,0.4575742) (-0.320000,0.4663560) (-0.306667,0.4749044) (-0.293333,0.4832051) (-0.280000,0.4912439) (-0.266667,0.4990075) (-0.253333,0.5064825) (-0.240000,0.5136563) (-0.226667,0.5205164) (-0.213333,0.5270512) (-0.200000,0.5332493) (-0.186667,0.5391000) (-0.173333,0.5445931) (-0.160000,0.5497191) (-0.146667,0.5544689) (-0.133333,0.5588342) (-0.120000,0.5628075) (-0.106667,0.5663815) (-0.093333,0.5695502) (-0.080000,0.5723078) (-0.066667,0.5746495) (-0.053333,0.5765711) (-0.040000,0.5780692) (-0.026667,0.5791412) (-0.013333,0.5797852) (0.000000,0.5800000) (0.013333,0.5797852) (0.026667,0.5791412) (0.040000,0.5780692) (0.053333,0.5765711) (0.066667,0.5746495) (0.080000,0.5723078) (0.093333,0.5695502) (0.106667,0.5663815) (0.120000,0.5628075) (0.133333,0.5588342) (0.146667,0.5544689) (0.160000,0.5497191) (0.173333,0.5445931) (0.186667,0.5391000) (0.200000,0.5332493) (0.213333,0.5270512) (0.226667,0.5205164) (0.240000,0.5136563) (0.253333,0.5064825) (0.266667,0.4990075) (0.280000,0.4912439) (0.293333,0.4832051) (0.306667,0.4749044) (0.320000,0.4663560) (0.333333,0.4575742) (0.346667,0.4485734) (0.360000,0.4393687) (0.373333,0.4299750) (0.386667,0.4204078) (0.400000,0.4106826) (0.413333,0.4008149) (0.426667,0.3908205) (0.440000,0.3807151) (0.453333,0.3705147) (0.466667,0.3602350) (0.480000,0.3498918) (0.493333,0.3395008) (0.506667,0.3290776) (0.520000,0.3186377) (0.533333,0.3081964) (0.546667,0.2977687) (0.560000,0.2873696) (0.573333,0.2770136) (0.586667,0.2667151) (0.600000,0.2564879) (0.613333,0.2463458) (0.626667,0.2363020) (0.640000,0.2263693) (0.653333,0.2165603) (0.666667,0.2068868) (0.680000,0.1973604) (0.693333,0.1879921) (0.706667,0.1787925) (0.720000,0.1697716) (0.733333,0.1609388) (0.746667,0.1523031) (0.760000,0.1438728) (0.773333,0.1356556) (0.786667,0.1276589) (0.800000,0.1198892) (0.813333,0.1123524) (0.826667,0.1050541) (0.840000,0.0979990) (0.853333,0.0911913) (0.866667,0.0846346) (0.880000,0.0783319) (0.893333,0.0722856) (0.906667,0.0664974) (0.920000,0.0609687) (0.933333,0.0556999) (0.946667,0.0506911) (0.960000,0.0459419) (0.973333,0.0414511) (0.986667,0.0372172) (1.000000,0.0332381) (1.013333,0.0295110) (1.026667,0.0260330) (1.040000,0.0228003) (1.053333,0.0198090) (1.066667,0.0170546) (1.080000,0.0145321) (1.093333,0.0122363) (1.106667,0.0101615) (1.120000,0.0083018) (1.133333,0.0066507) (1.146667,0.0052016) (1.160000,0.0039476) (1.173333,0.0028816) (1.186667,0.0019960) (1.200000,0.0012834) (1.213333,0.0007359) (1.226667,0.0003456) (1.240000,0.0001043) (1.253333,0.0000040) (1.266667,0.0000364) (1.280000,0.0001930) (1.293333,0.0004656) (1.306667,0.0008458) (1.320000,0.0013253) (1.333333,0.0018957) (1.346667,0.0025488) (1.360000,0.0032764) (1.373333,0.0040704) (1.386667,0.0049228) (1.400000,0.0058260) (1.413333,0.0067721) (1.426667,0.0077538) (1.440000,0.0087637) (1.453333,0.0097949) (1.466667,0.0108403) (1.480000,0.0118934) (1.493333,0.0129480) (1.506667,0.0139978) (1.520000,0.0150370) (1.533333,0.0160602) (1.546667,0.0170621) (1.560000,0.0180377) (1.573333,0.0189824) (1.586667,0.0198919) (1.600000,0.0207622) (1.613333,0.0215896) (1.626667,0.0223707) (1.640000,0.0231026) (1.653333,0.0237825) (1.666667,0.0244080) (1.680000,0.0249770) (1.693333,0.0254878) (1.706667,0.0259389) (1.720000,0.0263293) (1.733333,0.0266581) (1.746667,0.0269247) (1.760000,0.0271290) (1.773333,0.0272709) (1.786667,0.0273508) (1.800000,0.0273693) (1.813333,0.0273271) (1.826667,0.0272254) (1.840000,0.0270654) (1.853333,0.0268487) (1.866667,0.0265770) (1.880000,0.0262522) (1.893333,0.0258764) (1.906667,0.0254518) (1.920000,0.0249809) (1.933333,0.0244662) (1.946667,0.0239104) (1.960000,0.0233163) (1.973333,0.0226867) (1.986667,0.0220247) (2.000000,0.0213332) (2.013333,0.0206154) (2.026667,0.0198744) (2.040000,0.0191133) (2.053333,0.0183353) (2.066667,0.0175436) (2.080000,0.0167413) (2.093333,0.0159316) (2.106667,0.0151175) (2.120000,0.0143022) (2.133333,0.0134885) (2.146667,0.0126794) (2.160000,0.0118778) (2.173333,0.0110863) (2.186667,0.0103076) (2.200000,0.0095443) (2.213333,0.0087988) (2.226667,0.0080734) (2.240000,0.0073702) (2.253333,0.0066912) (2.266667,0.0060385) (2.280000,0.0054136) (2.293333,0.0048182) (2.306667,0.0042538) (2.320000,0.0037216) (2.333333,0.0032228) (2.346667,0.0027584) (2.360000,0.0023291) (2.373333,0.0019356) (2.386667,0.0015784) (2.400000,0.0012578) (2.413333,0.0009741) (2.426667,0.0007273) (2.440000,0.0005172) (2.453333,0.0003436) (2.466667,0.0002061) (2.480000,0.0001042) (2.493333,0.0000371) (2.506667,0.0000040) (2.520000,0.0000041) (2.533333,0.0000363) (2.546667,0.0000995) (2.560000,0.0001923) (2.573333,0.0003136) (2.586667,0.0004617) (2.600000,0.0006353) (2.613333,0.0008327) (2.626667,0.0010523) (2.640000,0.0012923) (2.653333,0.0015511) (2.666667,0.0018269) (2.680000,0.0021177) (2.693333,0.0024218) (2.706667,0.0027374) (2.720000,0.0030624) (2.733333,0.0033951) (2.746667,0.0037336) (2.760000,0.0040761) (2.773333,0.0044207) (2.786667,0.0047656) (2.800000,0.0051091)};
\draw[very thick]plot coordinates{(-2.800000,0.0116154) (-2.786667,0.0114606) (-2.773333,0.0112039) (-2.760000,0.0108491) (-2.746667,0.0104015) (-2.733333,0.0098685) (-2.720000,0.0092589) (-2.706667,0.0085833) (-2.693333,0.0078534) (-2.680000,0.0070824) (-2.666667,0.0062845) (-2.653333,0.0054744) (-2.640000,0.0046676) (-2.626667,0.0038798) (-2.613333,0.0031266) (-2.600000,0.0024235) (-2.586667,0.0017854) (-2.573333,0.0012262) (-2.560000,0.0007589) (-2.546667,0.0003952) (-2.533333,0.0001450) (-2.520000,0.0000165) (-2.506667,0.0000161) (-2.493333,0.0001479) (-2.480000,0.0004138) (-2.466667,0.0008134) (-2.453333,0.0013439) (-2.440000,0.0020003) (-2.426667,0.0027750) (-2.413333,0.0036583) (-2.400000,0.0046386) (-2.386667,0.0057018) (-2.373333,0.0068327) (-2.360000,0.0080140) (-2.346667,0.0092275) (-2.333333,0.0104539) (-2.320000,0.0116732) (-2.306667,0.0128654) (-2.293333,0.0140101) (-2.280000,0.0150876) (-2.266667,0.0160789) (-2.253333,0.0169661) (-2.240000,0.0177327) (-2.226667,0.0183641) (-2.213333,0.0188477) (-2.200000,0.0191732) (-2.186667,0.0193330) (-2.173333,0.0193223) (-2.160000,0.0191391) (-2.146667,0.0187847) (-2.133333,0.0182632) (-2.120000,0.0175820) (-2.106667,0.0167516) (-2.093333,0.0157853) (-2.080000,0.0146994) (-2.066667,0.0135126) (-2.053333,0.0122460) (-2.040000,0.0109227) (-2.026667,0.0095674) (-2.013333,0.0082063) (-2.000000,0.0068662) (-1.986667,0.0055745) (-1.973333,0.0043583) (-1.960000,0.0032444) (-1.946667,0.0022584) (-1.933333,0.0014246) (-1.920000,0.0007650) (-1.906667,0.0002996) (-1.893333,0.0000454) (-1.880000,0.0000161) (-1.866667,0.0002221) (-1.853333,0.0006698) (-1.840000,0.0013617) (-1.826667,0.0022962) (-1.813333,0.0034672) (-1.800000,0.0048646) (-1.786667,0.0064738) (-1.773333,0.0082763) (-1.760000,0.0102497) (-1.746667,0.0123678) (-1.733333,0.0146014) (-1.720000,0.0169182) (-1.706667,0.0192836) (-1.693333,0.0216612) (-1.680000,0.0240134) (-1.666667,0.0263018) (-1.653333,0.0284881) (-1.640000,0.0305345) (-1.626667,0.0324048) (-1.613333,0.0340647) (-1.600000,0.0354826) (-1.586667,0.0366302) (-1.573333,0.0374832) (-1.560000,0.0380218) (-1.546667,0.0382312) (-1.533333,0.0381019) (-1.520000,0.0376305) (-1.506667,0.0368192) (-1.493333,0.0356769) (-1.480000,0.0342186) (-1.466667,0.0324654) (-1.453333,0.0304449) (-1.440000,0.0281904) (-1.426667,0.0257407) (-1.413333,0.0231399) (-1.400000,0.0204364) (-1.386667,0.0176828) (-1.373333,0.0149346) (-1.360000,0.0122496) (-1.346667,0.0096873) (-1.333333,0.0073074) (-1.320000,0.0051693) (-1.306667,0.0033307) (-1.293333,0.0018469) (-1.280000,0.0007694) (-1.266667,0.0001453) (-1.253333,0.0000161) (-1.240000,0.0004167) (-1.226667,0.0013746) (-1.213333,0.0029092) (-1.200000,0.0050314) (-1.186667,0.0077423) (-1.173333,0.0110335) (-1.160000,0.0148866) (-1.146667,0.0192729) (-1.133333,0.0241538) (-1.120000,0.0294807) (-1.106667,0.0351953) (-1.093333,0.0412306) (-1.080000,0.0475111) (-1.066667,0.0539540) (-1.053333,0.0604702) (-1.040000,0.0669653) (-1.026667,0.0733413) (-1.013333,0.0794976) (-1.000000,0.0853329) (-0.986667,0.0907469) (-0.973333,0.0956416) (-0.960000,0.0999235) (-0.946667,0.1035055) (-0.933333,0.1063081) (-0.920000,0.1082617) (-0.906667,0.1093085) (-0.893333,0.1094032) (-0.880000,0.1085159) (-0.866667,0.1066323) (-0.853333,0.1037557) (-0.840000,0.0999079) (-0.826667,0.0951299) (-0.813333,0.0894830) (-0.800000,0.0830488) (-0.786667,0.0759295) (-0.773333,0.0682483) (-0.760000,0.0601481) (-0.746667,0.0517919) (-0.733333,0.0433611) (-0.720000,0.0350550) (-0.706667,0.0270885) (-0.693333,0.0196914) (-0.680000,0.0131054) (-0.666667,0.0075828) (-0.653333,0.0033833) (-0.640000,0.0007720) (-0.626667,0.0000161) (-0.613333,0.0013823) (-0.600000,0.0051336) (-0.586667,0.0115262) (-0.573333,0.0208063) (-0.560000,0.0332071) (-0.546667,0.0489453) (-0.533333,0.0682183) (-0.520000,0.0912013) (-0.506667,0.1180441) (-0.493333,0.1488689) (-0.480000,0.1837676) (-0.466667,0.2227995) (-0.453333,0.2659898) (-0.440000,0.3133277) (-0.426667,0.3647653) (-0.413333,0.4202165) (-0.400000,0.4795567) (-0.386667,0.5426225) (-0.373333,0.6092123) (-0.360000,0.6790863) (-0.346667,0.7519684) (-0.333333,0.8275471) (-0.320000,0.9054773) (-0.306667,0.9853832) (-0.293333,1.0668603) (-0.280000,1.1494785) (-0.266667,1.2327856) (-0.253333,1.3163104) (-0.240000,1.3995672) (-0.226667,1.4820588) (-0.213333,1.5632819) (-0.200000,1.6427303) (-0.186667,1.7199001) (-0.173333,1.7942936) (-0.160000,1.8654241) (-0.146667,1.9328202) (-0.133333,1.9960299) (-0.120000,2.0546250) (-0.106667,2.1082048) (-0.093333,2.1564001) (-0.080000,2.1988764) (-0.066667,2.2353370) (-0.053333,2.2655262) (-0.040000,2.2892312) (-0.026667,2.3062844) (-0.013333,2.3165650) (0.000000,2.3200000) (0.013333,2.3165650) (0.026667,2.3062844) (0.040000,2.2892312) (0.053333,2.2655262) (0.066667,2.2353370) (0.080000,2.1988764) (0.093333,2.1564001) (0.106667,2.1082048) (0.120000,2.0546250) (0.133333,1.9960299) (0.146667,1.9328202) (0.160000,1.8654241) (0.173333,1.7942936) (0.186667,1.7199001) (0.200000,1.6427303) (0.213333,1.5632819) (0.226667,1.4820588) (0.240000,1.3995672) (0.253333,1.3163104) (0.266667,1.2327856) (0.280000,1.1494785) (0.293333,1.0668603) (0.306667,0.9853832) (0.320000,0.9054773) (0.333333,0.8275471) (0.346667,0.7519684) (0.360000,0.6790863) (0.373333,0.6092123) (0.386667,0.5426225) (0.400000,0.4795567) (0.413333,0.4202165) (0.426667,0.3647653) (0.440000,0.3133277) (0.453333,0.2659898) (0.466667,0.2227995) (0.480000,0.1837676) (0.493333,0.1488689) (0.506667,0.1180441) (0.520000,0.0912013) (0.533333,0.0682183) (0.546667,0.0489453) (0.560000,0.0332071) (0.573333,0.0208063) (0.586667,0.0115262) (0.600000,0.0051336) (0.613333,0.0013823) (0.626667,0.0000161) (0.640000,0.0007720) (0.653333,0.0033833) (0.666667,0.0075828) (0.680000,0.0131054) (0.693333,0.0196914) (0.706667,0.0270885) (0.720000,0.0350550) (0.733333,0.0433611) (0.746667,0.0517919) (0.760000,0.0601481) (0.773333,0.0682483) (0.786667,0.0759295) (0.800000,0.0830488) (0.813333,0.0894830) (0.826667,0.0951299) (0.840000,0.0999079) (0.853333,0.1037557) (0.866667,0.1066323) (0.880000,0.1085159) (0.893333,0.1094032) (0.906667,0.1093085) (0.920000,0.1082617) (0.933333,0.1063081) (0.946667,0.1035055) (0.960000,0.0999235) (0.973333,0.0956416) (0.986667,0.0907469) (1.000000,0.0853329) (1.013333,0.0794976) (1.026667,0.0733413) (1.040000,0.0669653) (1.053333,0.0604702) (1.066667,0.0539540) (1.080000,0.0475111) (1.093333,0.0412306) (1.106667,0.0351953) (1.120000,0.0294807) (1.133333,0.0241538) (1.146667,0.0192729) (1.160000,0.0148866) (1.173333,0.0110335) (1.186667,0.0077423) (1.200000,0.0050314) (1.213333,0.0029092) (1.226667,0.0013746) (1.240000,0.0004167) (1.253333,0.0000161) (1.266667,0.0001453) (1.280000,0.0007694) (1.293333,0.0018469) (1.306667,0.0033307) (1.320000,0.0051693) (1.333333,0.0073074) (1.346667,0.0096873) (1.360000,0.0122496) (1.373333,0.0149346) (1.386667,0.0176828) (1.400000,0.0204364) (1.413333,0.0231399) (1.426667,0.0257407) (1.440000,0.0281904) (1.453333,0.0304449) (1.466667,0.0324654) (1.480000,0.0342186) (1.493333,0.0356769) (1.506667,0.0368192) (1.520000,0.0376305) (1.533333,0.0381019) (1.546667,0.0382312) (1.560000,0.0380218) (1.573333,0.0374832) (1.586667,0.0366302) (1.600000,0.0354826) (1.613333,0.0340647) (1.626667,0.0324048) (1.640000,0.0305345) (1.653333,0.0284881) (1.666667,0.0263018) (1.680000,0.0240134) (1.693333,0.0216612) (1.706667,0.0192836) (1.720000,0.0169182) (1.733333,0.0146014) (1.746667,0.0123678) (1.760000,0.0102497) (1.773333,0.0082763) (1.786667,0.0064738) (1.800000,0.0048646) (1.813333,0.0034672) (1.826667,0.0022962) (1.840000,0.0013617) (1.853333,0.0006698) (1.866667,0.0002221) (1.880000,0.0000161) (1.893333,0.0000454) (1.906667,0.0002996) (1.920000,0.0007650) (1.933333,0.0014246) (1.946667,0.0022584) (1.960000,0.0032444) (1.973333,0.0043583) (1.986667,0.0055745) (2.000000,0.0068662) (2.013333,0.0082063) (2.026667,0.0095674) (2.040000,0.0109227) (2.053333,0.0122460) (2.066667,0.0135126) (2.080000,0.0146994) (2.093333,0.0157853) (2.106667,0.0167516) (2.120000,0.0175820) (2.133333,0.0182632) (2.146667,0.0187847) (2.160000,0.0191391) (2.173333,0.0193223) (2.186667,0.0193330) (2.200000,0.0191732) (2.213333,0.0188477) (2.226667,0.0183641) (2.240000,0.0177327) (2.253333,0.0169661) (2.266667,0.0160789) (2.280000,0.0150876) (2.293333,0.0140101) (2.306667,0.0128654) (2.320000,0.0116732) (2.333333,0.0104539) (2.346667,0.0092275) (2.360000,0.0080140) (2.373333,0.0068327) (2.386667,0.0057018) (2.400000,0.0046386) (2.413333,0.0036583) (2.426667,0.0027750) (2.440000,0.0020003) (2.453333,0.0013439) (2.466667,0.0008134) (2.480000,0.0004138) (2.493333,0.0001479) (2.506667,0.0000161) (2.520000,0.0000165) (2.533333,0.0001450) (2.546667,0.0003952) (2.560000,0.0007589) (2.573333,0.0012262) (2.586667,0.0017854) (2.600000,0.0024235) (2.613333,0.0031266) (2.626667,0.0038798) (2.640000,0.0046676) (2.653333,0.0054744) (2.666667,0.0062845) (2.680000,0.0070824) (2.693333,0.0078534) (2.706667,0.0085833) (2.720000,0.0092589) (2.733333,0.0098685) (2.746667,0.0104015) (2.760000,0.0108491) (2.773333,0.0112039) (2.786667,0.0114606) (2.800000,0.0116154)};
\node[left]at(0,.58){$1$};\node[left]at(0,2.32){$4$};
\node[right]at(1.4,1.5){细线：$a$};\node[right]at(1.4,1.1){粗线：$2a$};
\node[below]at(0,-.5){$u=\pi ax/(\lambda L)$};
}
图\ref{fig:o3energya}比较同一照明强度下的两个孔宽。屏上近轴强度写成$I(x)=Ca^2\operatorname{sinc}^2[\pi ax/(\lambda L)]$，于是
\begin{align*}
 \int I(x)\dd x
 &=Ca^2\frac{\lambda L}{\pi a}
 \int\operatorname{sinc}^2u\dd u\\
 &=C\lambda La.
\end{align*}
因此$a\to2a$时峰高乘4、横向尺度减半、积分功率乘2。图中两条曲线按sinc平方绘制，旁瓣也参与上述全平面积分。对有限高度的矩形孔，还要沿另一屏坐标积分，最终功率正比真实孔面积。此守恒关系以同一入射场强、同一无损传播条件为前提；若固定的是总入射功率，孔面振幅就应随孔宽重新归一化。
\end{solution}

\begin{exercise}\label{ex:o3oblique}
周期$b$的透射光栅受入射角$i$的平面波照射，入射和出射横向波矢分别为$k\sin i$、$k\sin\theta$。推导光栅方程和同级角色散，并讨论入射角是否直接改变$mN$。
\Diagram{fig:o3obliqueq}{斜入射与出射的横向波矢投影}{\draw[curve](-2,0)--(2,0);\draw[guide](0,-1.7)--(0,2.2);\draw[->](-1.2,-1.7)--(0,0);\draw[->](0,0)--(1.7,2.1);\node[left]at(-.6,-1){$i$};\node[right]at(.9,1.1){$\theta$};}
\end{exercise}
\begin{solution}
图\ref{fig:o3obliqueq}沿孔径横坐标相邻周期的净相位差为$kb(\sin i-\sin\theta)$。相长条件取整数级次约定为$b(\sin\theta-\sin i)=m\lambda$。固定$i$，微分得到$\dd\theta/\dd\lambda=m/(b\cos\theta)$。
\Diagram{fig:o3obliquea}{首零角宽与谱线位移共享同一投影因子}{\draw[axis](0,0)--(4.6,0)node[right]{$\theta$};\draw[curve](1,0)..controls(1.6,2.6)and(1.8,2.6)..(2.4,0);\draw[guide](2,0)..controls(2.6,2.6)and(2.8,2.6)..(3.4,0);\node[below]at(2.3,-.3){$\Delta\theta$};}
首零半宽仍为$\lambda/(Nb\cos\theta)$，两式之比使投影因子消去，故按瑞利判据仍有$\lambda/\Delta\lambda=mN$。图\ref{fig:o3obliquea}表示在固定级次、固定实际被照亮缝数时的比较。斜入射会改变足迹宽度、可传播级次和效率，所以在固定入射束截面条件下，$N$本身可能改变。
\end{solution}

\begin{exercise}\label{ex:o3defocus}
在夫琅禾费装置中，观察屏偏离透镜后焦面。为什么衍射场重新出现孔径二次相位？以一维自由传播说明保留哪些项。
\Diagram{fig:o3defocusq}{焦面前后不同观察平面}{\OptAxis\draw[<->](1.5,-1.5)--(1.5,1.5);\draw[guide](4,-1.7)--(4,1.7);\draw[curve](5,-1.7)--(5,1.7);\node[above]at(4,1.7){焦面};\node[above]at(5,2.2){偏移屏};}
\end{exercise}
\begin{solution}
图\ref{fig:o3defocusq}中，令透镜前孔面场为$U_A(\xi)$。薄透镜相位为$\exp[-ik\xi^2/(2f)]$，向屏传播$L$的近轴核为$\exp[ik(x-\xi)^2/(2L)]$，两者相乘并展开平方，得到
\[
 U(x,L)=C_L\ee^{ikx^2/(2L)}\int U_A(\xi)
 \ee^{\frac{ik\xi^2}{2}(1/L-1/f)}\ee^{-ikx\xi/L}\dd\xi.
\]
观察点公共相位$\exp[ikx^2/(2L)]$不改变单点强度，孔径内的二次相位却改变各面元的相干叠加，不能略去。
\Diagram{fig:o3defocusa}{只有$L=f$时二次相位恰好抵消}{\begin{scope}[yscale=.85]\draw[axis](-2.2,0)--(2.4,0)node[right]{$\xi$};\draw[axis](0,-1.5)--(0,2)node[above]{残余相位};\draw[curve,domain=-1.8:1.8,samples=60]plot(\x,{.4*\x*\x});\draw[guide,domain=-1.8:1.8,samples=60]plot(\x,{-.3*\x*\x});\node[right]at(1.8,1.3){$L<f$};\node[right]at(1.8,-1){$L>f$};\end{scope}}
图\ref{fig:o3defocusa}中$L=f$时二次项准确抵消，远场成为孔函数在$q=kx/f$处的傅里叶变换。若$L=f+\Delta L$，则$1/L-1/f=-\Delta L/[f(f+\Delta L)]\simeq-\Delta L/f^2$。以孔径全宽$a$、边缘$\xi=a/2$计，最大残余相位大小近似为
\[
 |\Delta\phi_{\rm edge}|\simeq\frac{k a^2|\Delta L|}{8f^2}.
\]
要求它远小于1，得到$|\Delta L|\ll8f^2/(ka^2)$。若改用半宽$a$记号，分母系数随之改变，不能混用。由此可见，判断离焦是否可略，须同时考虑波长、孔径与焦距；单独的$|\Delta L|/f\ll1$只保证几何距离接近，并不保证孔内相位误差小。
\end{solution}

\begin{exercise}\label{ex:o3aperturephase}
两条相同狭缝中心间隔$b$，在一条缝前加入均匀相位片$\phi_0$。求衍射分布；若$\phi_0=\pi$，轴上是否仍有中央主极大？
\Diagram{fig:o3phaseq}{相位片只改变一缝整体相位，不改变其宽度}{\begin{scope}[yscale=.75]\draw[curve](0,-1.2)rectangle(.4,-.4);\draw[curve](0,.4)rectangle(.4,1.2);\draw[very thick](-.2,.4)--(-.2,1.2);\node[left]at(-.2,.8){$\phi_0$};\draw[axis](.8,0)--(4,0)node[right]{远场};\end{scope}}
\end{exercise}
\begin{solution}
取图\ref{fig:o3phaseq}两缝中心坐标为0和$b$，则$U=U_a[1+\ee^{i(\phi_0-kb\sin\theta)}]$，故$I=4I_a(0)\operatorname{sinc}^2\beta\cos^2[(\phi_0-kb\sin\theta)/2]$。
\Diagram{fig:o3phasea}{附加$\pi$相位把轴上两缝相量变为反向}{\draw[->,thick](0,0)--(2,0)node[right]{$A$};\draw[->,thick](0,0)--(-2,0)node[left]{$-A$};\node[above]at(0,.35){$\theta=0$};}
在$\phi_0=\pi$时，图\ref{fig:o3phasea}轴上相消，强度为零。单缝包络仍以轴为中心，但干涉细纹平移半个周期；不能把包络中心自动称为总强度的主极大。若相位片有吸收，振幅失配会留下轴上残光。
\end{solution}

